\documentclass[11pt]{amsart}
\usepackage[T1]{fontenc}
\usepackage{lmodern,amsmath,amssymb,amsthm,mathtools,microtype,booktabs,array,longtable}
\usepackage[margin=1.06in]{geometry}
\usepackage[hidelinks]{hyperref}
\emergencystretch=2em
\numberwithin{equation}{section}
\newtheorem{theorem}{Theorem}[section]
\newtheorem{lemma}[theorem]{Lemma}
\newtheorem{proposition}[theorem]{Proposition}
\newtheorem{corollary}[theorem]{Corollary}
\theoremstyle{definition}\newtheorem{definition}[theorem]{Definition}
\theoremstyle{remark}\newtheorem{remark}[theorem]{Remark}
\newcommand{\Q}{\mathbb Q}\newcommand{\R}{\mathbb R}\newcommand{\C}{\mathbb C}\newcommand{\N}{\mathbb N}\newcommand{\Z}{\mathbb Z}\newcommand{\E}{\mathbb E}
\newcommand{\ip}[2]{\langle #1,#2\rangle}\newcommand{\norm}[1]{\lVert #1\rVert}
\newcommand{\epsc}{\varepsilon_{\mathrm c}}
\DeclareMathOperator{\osc}{osc}\DeclareMathOperator{\spanop}{span}\DeclareMathOperator{\End}{End}\DeclareMathOperator{\diag}{diag}\DeclareMathOperator{\tr}{tr}\DeclareMathOperator{\rad}{rad}\DeclareMathOperator{\TV}{TV}
\title[Finite actions and observation processes]{Finite Physical Actions\\and Repeatable Observation Processes}
\author{Qian Qi}
\date{27 September 2026}
\subjclass[2020]{Primary 93B15, 22C05; Secondary 68Q45, 52A27, 60J10}
\begin{document}
\begin{abstract}
We characterize finite-label causal realizations of compact-group
experiments with repeatable observations. For every fixed error below
one in total variation on the whole adaptive transcript, bounded
horizon-dependent width is equivalent to finiteness of the future
response quotient and to existence of an exact stationary finite
physical action. Below error one half, its cardinality is the eventual
minimum width. Arbitrary intervening register sizes are allowed in the
converse. Repeatable probes supply physical holds; no return condition
on the command alphabet is needed for this process theorem. Separately,
for terminal convex-valued outputs, eventual approximate physical
returns imply same-width, same-error stationarization. Compact
stochastic-semigroup purification then yields a finite physical action
with randomized initialization and an equality-inclusive quotient-radius
criterion. We retain the component, phase, effectivity and finite-group
complexity consequences. The process theorem is proved by common
statistical tests across a causal register, not by extending terminal
purification to path laws.
\end{abstract}
\maketitle
\noindent The resource counts persistent hidden labels, including unused
padding. The horizon, clock, real table descriptions, table construction,
exact arithmetic and exact sampling are uncharged. A numerical terminal
decoder is not a sampled answer register; a process error is one budget
for the whole control--output transcript.

\section{Introduction}
There are two distinct realization questions. A terminal numerical
interface asks for a correct mean decoder after a word. A causal
observation experiment asks for one consistent law of all outputs,
including controls selected from earlier outputs. With repeatable
probes, Section~\ref{sec:causal59} proves that bounded width at any
fixed total transcript error below one requires a finite future-response
quotient, and that this quotient itself gives an exact realization.
This conclusion is stronger than a radius threshold for numerical
marginals and has a different proof. The terminal theory is developed
in full thereafter.

We ask whether a free clock and a separate stochastic realization at
every horizon can reduce the limiting number of hidden labels needed to
preserve a continuous numerical interface. Under eventual approximate
returns, they cannot: the least clocked width converges to the stationary
minimum at the same closed error, and finite limits are eventually
attained. More strongly, arbitrarily many cuts of a fixed width already
force a stationary realization of that width, even when all other
registers are unrestricted.

For compact physical groups, stationary transition randomness and hidden
permutations above the physical identity can be removed without losing
a label or increasing error. Randomized initialization remains. This
identifies the stationary minimum with a finite continuous physical
action and yields an equality-inclusive finite-quotient criterion,
including groups with infinitely many components. A change of frame
identifies each phase minimum without a phase-register overhead.

The command-width resource immediately below excludes any sampled answer
register. The full proofs of the structural reductions, component-radius
localization, phase and profinite refinements, algebraic effectivity, and
explicit finite-group complexity are contained in this article. The
separate quantitative paper concerns orbit geometry and an explicit
rational exact--noisy separation; none of its theorems is a premise here.
\subsection{The interface and its two resources}
Unless a section explicitly specifies a different physical space, $H$ is a compact Hausdorff group and $\mathcal A$ is a finite nonempty alphabet whose letters $u_a\in H$ generate a dense subgroup. The execution convention is
\[
                  u_{a_1\cdots a_n}=u_{a_n}\cdots u_{a_1}.
\]
There are finite nonempty seed and query sets $\mathcal S,\mathcal J$ and continuous maps $F_{s,j}:H\to Y_j$. In the structural results $Y_j$ is a nonempty norm-compact convex subset of a real Banach space. The localization and algebraic sections specify finite-dimensional output spaces. For a categorical law, $Y_j=\Delta_{M_j}$ and the norm on differences is $\frac12\|\cdot\|_1$, namely total variation. A prescribed joint law is one query; incompatible measurements are not assigned an artificial joint law.

\begin{definition}[Stationary realization]\label{def:stationary55}
A realization on $k$ labels consists of initialization rows $\alpha_s\in\Delta_k$, row-stochastic matrices $T_a$ of order $k$, and decoder values $D_j(i)\in Y_j$, all independent of time and word length. Its error is
\[
       \sup_{s,j,w\in\mathcal A^*}
        \|\alpha_sT_{a_1}\cdots T_{a_n}D_j-F_{s,j}(u_w)\|_j.
\]
The empty word is included. Let $M_\varepsilon$ be the least $k$ at error at most $\varepsilon$, or $\infty$ when none exists. A permutation realization has permutation command matrices; its initialization need not be deterministic.
\end{definition}
The clocked resource $W_{N,\varepsilon}$ instead minimizes the maximum size of the registers $S_0,\ldots,S_N$ at a fixed horizon, with arbitrary row-stochastic transitions depending on time and horizon, and a terminal query decoder. Definition~\ref{def:machine54} states the model formally. Every advertised persistent label counts, including unused padding. Private information retained for future commands is part of the current label. The original width model does not charge the external clock, real table descriptions, their construction, exact arithmetic, or exact atomic sampling. Section~\ref{sec:complexity56} gives separate finite-description and fair-random-bit bounds; it does not reinterpret those resources as free computational memory. A numerical decoder vector is not an additional sampled-answer register.

For an open normal subgroup $U\le H$ put
\begin{equation}\label{eq:quotientradius55}
 r(U)=\max_{s,j,h\in H}\rad_{Y_j}(F_{s,j}(hU)),
 \qquad\rad_Y(A)=\min_{y\in Y}\sup_{x\in A}\|x-y\|.
\end{equation}
The coset images are compact by continuity. The radius depends continuously on the translate, and compactness gives the displayed extrema. The restriction of the center to the legal output set is essential. The component radius $R_0$ uses $U=H^\circ$, whether or not this subgroup is open; it is the infimum of the finite-quotient radii. At the critical error, the relevant issue is attainment of that infimum, not a limit from larger errors.


The principal dependencies are finite stationary witnesses and Ramsey
extraction for clock removal, classical recurrent-simplex structure for
purification, and physical-kernel averaging for descent. The associated
new quantifiers are the same width, the same closed numerical error,
randomized initialization and arbitrary intervening register widths.
The classical semigroup statements themselves are not claimed as new.

\section{Repeatable observations and causal realization}\label{sec:causal59}
Terminal numerical accuracy and accuracy of a generated observation process
are different requirements. Repeating a nondisturbing observation can
separate arbitrarily close, but distinct, numerical responses. We now
formulate this distinction within the same finite-label resource model.

Let $H$ be a compact Hausdorff group, let $\mathcal A$ be a finite nonempty
command alphabet generating a dense subgroup, and let $\mathcal S$ and
$\mathcal J$ be finite nonempty sets. For each $s,j$ fix a continuous law
\[
 p_{s,j}:H\longrightarrow\Delta_{M_j}.
\]
The physical state starts at $g=e$. A command $a$ replaces $g$ by $u_ag$
and emits a fixed null symbol. A probe $j$ leaves $g$ unchanged and emits
a fresh draw from $p_{s,j}(g)$, conditionally on the public past
and current physical state. At each step the next command or probe may be chosen from the entire
observed transcript by a randomized policy. The seed is specified at
initialization. Probes are repeatable statistical observations; when
$p_{s,j}$ is obtained from a density matrix, this is an oracle model, not
repeated nondisturbing measurement of one unknown quantum system.

\begin{definition}[Causal finite-label machine]\label{def:causal59}
At horizon $N$ a machine has registers $S_0,\ldots,S_N$, initialization
rows $E_s$, and, for each control $c$ and possible output $y$, a
nonnegative matrix $K_{t,c,y}:S_{t-1}\to S_t$ such that
$\sum_yK_{t,c,y}$ is row-stochastic. Its entries jointly specify the
emission and next label. The matrices may depend on $t,N$ but on the past
only through the retained label and current control. Its error is
\[
 \sup_{s,\pi}\TV\bigl(\mathbb P^{\pi}_{s,N},
                            \mathbb Q^{\pi}_{s,N}\bigr),
\]
where the laws include the complete control--output transcript under the
same policy $\pi$. Write $\mathsf{PW}_{N,\varepsilon}$ for the minimum
peak register size at error at most $\varepsilon$.
\end{definition}
Every advertised label and all retained private randomness are counted.
The horizon, external clock and arbitrary real tables remain free. The
policy is an external experimenter, not an uncharged private memory for
the machine. Outputs may influence future controls, but are not supplied
again as a free history input. A sampled output register, when charged,
adds the separate fixed output-alphabet cost. The error above is one
budget for the entire transcript, not a per-step error budget.

For $(s,g)$ define its future response function
\begin{equation}\label{eq:futuretype59}
 \Theta_{s,g}(x,j)=p_{s,j}(xg),\qquad x\in H,
 \qquad \mathcal X=\{\Theta_{s,g}:s\in\mathcal S,\ g\in H\}.
\end{equation}
Equip the finite product of continuous-function spaces with the uniform
norm. Put $n=|\mathcal X|$ when this set is finite, and $n=\infty$
otherwise. This is a behavioral quotient, not the dimension of a linear
predictive representation. In particular, infinitely many distinct
predictive rows do not by themselves imply infinite nonnegative
realization size for a general hidden Markov process
\cite[Theorem~2 and Section~4]{SJR2004}. Repeatable observations are the
additional input in the theorem below.

\begin{theorem}[Uniform process error and finite physical actions]
\label{thm:causalclassification59}
For every fixed $0\le\varepsilon<1$ the following are equivalent:
\begin{enumerate}
\item $\sup_N\mathsf{PW}_{N,\varepsilon}<\infty$;
\item $n<\infty$;
\item some open normal subgroup $U\lhd H$ makes every $p_{s,j}$ constant
      on each coset of $U$;
\item a stationary exact causal realization exists with permutation
      command updates and state-preserving probe emissions.
\end{enumerate}
If $n<\infty$ and $0\le\varepsilon<1/2$, then
\begin{equation}\label{eq:causalminimum59}
                   \mathsf{PW}_{N,\varepsilon}=n
                   \quad\hbox{for all sufficiently large }N.
\end{equation}
If $n=\infty$, then, for every $\varepsilon<1$ and every $k\ge1$,
there is $L_{k,\varepsilon}<\infty$ such that every admissible realization
satisfies
\begin{equation}\label{eq:causaloccupation59}
 \#\{t\in\{1,\ldots,N\}:|S_t|\le k\}\le L_{k,\varepsilon}.
\end{equation}
The same occupation conclusion holds for finite $n$, $\varepsilon<1/2$
and $k<n$. Intervening registers are unrestricted. No return condition
is imposed on the command alphabet: the probes themselves leave the
physical state unchanged and supply the padding used in the proof.
\end{theorem}
The threshold $1$ in the boundedness criterion is sharp: total variation
never exceeds one, so at tolerance one a single label always suffices.
The formula \eqref{eq:causalminimum59} is asserted only below $1/2$;
no unrestricted minimum formula for $1/2\le\varepsilon<1$ is implicit.

\begin{lemma}[A common separating itinerary]\label{lem:itinerary59}
Choose $m\ge2$ different functions in \eqref{eq:futuretype59}, represented
by executable seed--prefix pairs. For every $\delta>0$ there is one
finite deterministic control itinerary with disjoint output events
$A_1,\ldots,A_m$ such that its target law from the $i$th state assigns
probability at least $1-\delta$ to $A_i$.
\end{lemma}
\begin{proof}
Equality of two future response functions is preserved and reflected by
left multiplication of their physical states by any $u\in H$: replace
$x$ in \eqref{eq:futuretype59} by $xu$, a bijection of $H$. Positive
executable products are dense in $H$. Indeed positive powers approximate
the inverse of every generator in a compact group, so the closed positive
semigroup contains the dense generated group.

List the $\binom m2$ pairs. At each stage the current two states for that
pair still have different future response functions. By continuity and
positive-word density, a finite executable word followed by one probe
separates some coordinate of their output laws. Append that word and
record the coordinate, then proceed to the next pair. This constructs
one itinerary, not a family of incompatible pair-dependent suffixes.
No reset or inverse word is required. Let $J=\binom m2$, and let
$\mu_{i,r}$ be the selected coordinate probability at the $r$th probe
location when the initial state is the $i$th representative. There is
$\Delta>0$ such that every pair of rows of this finite matrix differs
by at least $\Delta$ in some coordinate.

At each location repeat the chosen probe $L$ times. These repetitions
leave the physical state fixed. The empirical coordinate frequencies
$\widehat\mu_r$ obey
\[
 \mathbb P_i\{\max_r|\widehat\mu_r-\mu_{i,r}|\ge\Delta/3\}
       \le 2J\exp(-2L\Delta^2/9).
\]
For completeness, the scalar bound follows from
$\E e^{z(B-\E B)}\le e^{z^2/8}$ for a Bernoulli variable, exponential
Markov inequality, and optimization in $z$; a union bound gives the
display. The moment bound follows by differentiating the logarithmic
moment generating function twice: its second derivative is the variance
of a tilted Bernoulli variable and is at most $1/4$.
Choose $L\ge1$ with the displayed upper bound at most $\delta$.
The events
$A_i=\{\max_r|\widehat\mu_r-\mu_{i,r}|<\Delta/3\}$ are disjoint,
because $2\Delta/3<\Delta$, and have the required probabilities.
\end{proof}

\begin{lemma}[Testing across a finite register]\label{lem:testingcut59}
Suppose $m$ equal-length seed--prefix pairs at cut $t$ admit an itinerary
of length $\ell$ as in Lemma~\ref{lem:itinerary59}. For every
$N\ge t+\ell$, an error-$\varepsilon$ causal realization satisfies
\begin{equation}\label{eq:testingcut59}
                  |S_t|\ge m(1-\varepsilon-\delta).
\end{equation}
If $\varepsilon+\delta<1/2$, the stronger bound $|S_t|\ge m$ holds.
\end{lemma}
\begin{proof}
Use the same deterministic suffix for all representatives and pad its
end with probes if necessary. Marginalize all prefix outputs. Write
$\alpha_i(z)$ for the resulting cut-label distribution and $Q_z$ for
the suffix-output law conditional on label $z$. Causality makes $Q_z$
independent of the chosen prefix. Total variation contracts under
marginalization, so correctness gives
\[
 1-\delta-\varepsilon
     \le\sum_z\alpha_i(z)Q_z(A_i).
\]
This uses no conditional correctness assertion after a rare observed
prefix. Summing in $i$, using $\alpha_i(z)\le1$, and then disjointness,
yields
\[
 m(1-\delta-\varepsilon)
 \le\sum_z\sum_iQ_z(A_i)\le |S_t|.
\]
If $1-\delta-\varepsilon>1/2$, each $i$ has a positive-mass label $z_i$
with $Q_{z_i}(A_i)>1/2$. Disjointness forces the labels $z_i$ to be
different. This proves the stronger assertion.
\end{proof}

\begin{proof}[Proof of Theorem~\ref{thm:causalclassification59}]
The map $g\mapsto\Theta_{s,g}$ is continuous in the uniform norm.
This follows from continuity on the compact product $H\times H$ and a
finite-cover argument, uniformly in the first coordinate. Thus each
seed's image is compact. Executable images are dense in it. If
$\mathcal X$ is infinite, arbitrarily many distinct types therefore
have executable representatives; if it is finite, every type has one.

Choose finitely many such representatives and pad their prefixes to a
common length $t_0$ by any fixed probe. Their physical states are
unchanged. Apply Lemma~\ref{lem:itinerary59}. For every
$t_0\le t\le N-\ell$, pad further with the same probe, marginalize its
outputs, and apply Lemma~\ref{lem:testingcut59}. The itinerary and
$\ell$ do not depend on $t,N$ or on the machine.

If $n=\infty$, choose $m$ with $m(1-\varepsilon)>k$ and then choose
$\delta>0$ with $m(1-\varepsilon-\delta)>k$. Every interior cut has
width greater than $k$. At most $t_0+\ell$ positive cuts lie outside
this interval, including when it is empty. This proves
\eqref{eq:causaloccupation59}, divergence, and the implication (i)$\Rightarrow$(ii).
For finite $n$ and $\varepsilon<1/2$, use $m=n$ and
$0<\delta<1/2-\varepsilon$ when $n\ge2$. All sufficiently interior
cuts require $n$ labels. The case $n=1$ follows from nonempty registers.

For finite $\mathcal X$, the maps
\[
 \Theta(x,j)\longmapsto\Theta(xu,j),\qquad u\in H,
\]
define a continuous permutation action of $H$ on $\mathcal X$.
Continuity follows from uniform continuity of each of the finitely many
functions. Its kernel $U$ is open and normal. Evaluating at $x=e$
shows that every $p_{s,j}$ is constant on $U$-cosets, proving
(ii)$\Rightarrow$(iii). Conversely (iii) gives at most
$|\mathcal S|[H:U]$ response types. More economically, use the $n$
distinct types as labels, initialize at $\Theta_{s,e}$, update by the
permutation above, and emit $\Theta(e,j)$ on probe $j$ without changing
the label. This produces the exact conditional response at each step.
Induction on transcript length proves equality for every adaptive
policy, hence (iv) and (i). It also supplies the upper bound $n$ in
\eqref{eq:causalminimum59}. The occupation assertion for $k<n$ follows
from the stronger testing bound just established.
\end{proof}

\begin{corollary}[Connected groups]\label{cor:connectedprocess59}
If $H$ is connected, uniformly bounded causal width at any fixed total
transcript error below one is possible exactly when every $p_{s,j}$ is
constant on $H$. If at least one is nonconstant, every fixed width has
the occupation bound \eqref{eq:causaloccupation59}.
\end{corollary}
\begin{proof}
A continuous action of a connected group on a finite discrete set is
trivial. The assertion follows from the theorem.
\end{proof}

\begin{remark}[Why terminal purification is insufficient]\label{rem:correlation59}
A two-label variable chosen once with equal masses and then repeatedly
reported has the correct Bernoulli$(1/2)$ marginal at every time, but its
length-$L$ law is supported on the two constant strings. Its total
variation distance from $L$ independent fair bits is $1-2^{1-L}$.
Thus even exact numerical marginals do not supply a path-law realization.
The causal theorem concerns an actual nonnegative emission instrument,
not a list of separately correct decoders. It neither classifies
arbitrary irreversible hidden-state systems nor asserts a measurement
procedure for noncommuting quantum observables.
\end{remark}

\section{Stationarization from recurrent narrow cuts}\label{sec:ramsey56}
The first reduction does not require reversible physical dynamics. In this section only, replace $H$ by a topological monoid $G$ with identity $1$, retain a finite command alphabet, and let $F_{s,j}:G\to Y_j$ be continuous. The output sets are norm-compact convex subsets of real Banach spaces. Define $M_\varepsilon$ and $W_{N,\varepsilon}$ by the same stochastic models. Compactness of the physical monoid is not required.

\begin{definition}[Eventual approximate returns]\label{def:returns56}
The alphabet has eventual approximate returns if, for every neighborhood $V$ of $1$, there is an integer $g(V)\ge0$ such that every integer $n\ge g(V)$ is the length of a word with product in $V$. Cofinite exact returns mean that a single $g$ works with product exactly $1$.
\end{definition}
The statement concerns physical products only. The hidden rows on a return word need not be close to an identity matrix. Approximation here is eventual at every length, not merely along a return subsequence.

\begin{theorem}[Width-preserving stationarization]\label{thm:stationarization56}
Assume eventual approximate returns. For every $\varepsilon\ge0$ and integer $k\ge1$, the following are equivalent:
\begin{enumerate}
\item a stationary realization with at most $k$ labels has error at most $\varepsilon$;
\item among error-$\varepsilon$ clocked realizations, the number of positive cuts having width at most $k$ is unbounded.
\end{enumerate}
Consequently, if $M_\varepsilon>k$, there is $L_{k,\varepsilon}<\infty$, independent of the horizon, such that every error-$\varepsilon$ realization satisfies
\begin{equation}\label{eq:occupation56}
 \#\{t\in\{1,\ldots,N\}:|S_t|\le k\}\le L_{k,\varepsilon}.
\end{equation}
All other registers may be arbitrarily large. Moreover,
\begin{equation}\label{eq:limitminimum56}
                  \lim_{N\to\infty}W_{N,\varepsilon}=M_\varepsilon
                  \quad\hbox{in }\N\cup\{\infty\}.
\end{equation}
If the right side is finite, equality holds for every sufficiently large $N$.
\end{theorem}
The theorem preserves the proposed width, including randomized initialization. It does not assert that the microscopic rows of the clocked machine become stationary. They can be entirely unrelated at different times.

\subsection{Finite witnesses and homogeneous transition tables}
For a matrix $A$ write $\|A\|_{\mathrm r}=\max_i\sum_j|A_{ij}|$. Put $a=|\mathcal A|$ and $D_* =\max(1,\max_j\operatorname{diam}Y_j)$. Let $\mathcal X_k$ be the compact parameter space of stationary $k$-label machines. For $m\ge0$ define
\[
 e_{k,m}=\min_{X\in\mathcal X_k}\max_{s,j,\,|w|\le m}
               \|X_{s,j}(w)-F_{s,j}(u_w)\|_j.
\]
The finite maximum is continuous, so the minimum is attained. If no stationary error-$\varepsilon$ machine exists, compactness of $\mathcal X_k$ and the finite intersection property give an $m$ with
\begin{equation}\label{eq:finitewitness56}
                       e_{k,m}=\varepsilon+\gamma,\qquad\gamma>0.
\end{equation}
Indeed the nested closed sets defined by the constraints $|w|\le m$ otherwise have a common point. This argument also applies when the full infinite-word error is not a continuous function of the parameters.

We shall use the finite edge-coloring theorem of Ramsey \cite{Ramsey}. Denote by $R_q(h)$ an integer such that every coloring of the edges of a complete graph on $R_q(h)$ vertices with at most $q$ colors has a monochromatic $h$-vertex subgraph. Only finiteness is essential.

\begin{remark}[An explicit but coarse Ramsey bound]
For definiteness one may take
\begin{equation}\label{eq:ramseybound56}
                   R_q(h)\le q^{qh+1}\qquad(q\ge2,\ h\ge2).
\end{equation}
To see this elementary bound, successively select a vertex and retain a largest color class among its remaining neighbors. With $q^{qh+1}$ initial vertices this produces at least $q(h-1)+1$ selected vertices, each of whose edges to later selected vertices has one color. At least $h$ of their associated colors agree; these vertices form the required subgraph. The deliberately excessive exponent avoids any endpoint rounding issue.
\end{remark}

Partition $[0,1]$ into at most $\lceil k/\delta\rceil+1$ intervals of length at most $\delta/k$. Two tuples in the same entrywise cell of
\[
                  \mathsf{Stoch}(k)^{\{0\}\cup\mathcal A}
\]
have corresponding matrices within $\delta$ in $\|\cdot\|_{\mathrm r}$. The number of occupied cells is at most
\begin{equation}\label{eq:colors56}
             q=\max\bigl(2,(\lceil k/\delta\rceil+1)^{(a+1)k^2}\bigr).
\end{equation}
Each representative used below is chosen from an occupied cell and is therefore stochastic; an entrywise rounded matrix need not be stochastic and is not used as a representative.

\begin{proof}[Proof of Theorem~\ref{thm:stationarization56}]
The forward implication follows by using one stationary machine at every horizon. For the converse, suppose \eqref{eq:finitewitness56}. Choose $0<\eta<\gamma/4$. Continuity of multiplication and of the finitely many $F_{s,j}$ gives a neighborhood $V$ of $1$ with the following property. In any word of length at most $m$, inserting at most $m+3$ factors from $V$, including before and after the word, changes every output by less than $\eta$. There are only finitely many relevant command words and insertion patterns; continuity at the tuple of identity factors, followed by a finite intersection of neighborhoods, proves the assertion. This uses neither an invariant metric nor commutativity.

Let $g=g(V)$, put $b=g+1$, and choose
\[
       0<\delta\le \min\bigl(1,\gamma/(4D_*(m+1))\bigr).
\]
Take a clocked realization with so many narrow cuts that, after discarding cuts $t<g$ or $t>N-g$ and greedily retaining cuts separated by at least $b$, at least $R_q(m+2)$ remain. Such a choice is possible by hypothesis; at most $2g$ cuts are discarded and each retained cut accounts for at most $b$ of the others.

Identify each retained register with a subset of $\{1,\ldots,k\}$. Extend rows from dummy labels arbitrarily to stochastic rows without altering any row on the actual register. For retained cuts $t_i<t_j$ write $\ell=t_j-t_i\ge g+1$. Choose a length-$\ell$ return word with product in $V$, and let $A_{ij}$ be its actual composite kernel between these cuts. For each letter $c\in\mathcal A$, choose the word consisting of $c$ followed by a length-$(\ell-1)$ return word with product in $V$, and let $B_{ij,c}$ be its composite kernel. Thus the physical products of these words are respectively in $V$ and of the form $v u_c$ with $v\in V$. All the matrices are of order $k$, however large the intermediate registers may be.

Color the edge $\{i,j\}$, oriented from the earlier cut to the later cut, by the cell of $(A_{ij},(B_{ij,c})_c)$. Ramsey's theorem gives $m+2$ ordered cuts
\[
                         t_0<t_1<\cdots<t_{m+1}
\]
whose edges have one color. Fix a stochastic representative $(E,(B_c)_c)$ from that color. Choose fixed return words for the prefix of length $t_0$ and the suffix of length $N-t_{m+1}$, both with products in $V$. Let $\alpha_s$ be the actual state distribution after the prefix. Let $D_j(i)$ be the conditional expected original decoder after the suffix, given label $i$ at its start, with arbitrary legal values on dummy labels. Convexity gives $D_j(i)\in Y_j$.

Consider the stationary machine
\begin{equation}\label{eq:extractedmachine56}
                  \alpha_s,\qquad T_c=B_c,\qquad d_j=ED_j.
\end{equation}
Its decoder is legal. To test $w=c_1\cdots c_\ell$, $0\le\ell\le m$, use the first $\ell$ consecutive clique edges with kernels $B_{i-1,i,c_i}$ and then the single edge from $t_\ell$ to $t_{m+1}$ with kernel $A_{\ell,m+1}$. For $\ell=0$ there is only this last edge. Add the fixed prefix and suffix. This is a genuine word of length $N$; its physical product is obtained from $u_w$ by at most $\ell+3$ insertions from $V$. Its target output is consequently within $\eta$ of $F_{s,j}(u_w)$.

Each of its $\ell+1$ edge matrices differs from the corresponding matrix in \eqref{eq:extractedmachine56} by at most $\delta$. Stochastic multiplication contracts the $\ell^1$ distance between probability rows. A telescoping expansion, and subtraction of any fixed point of $Y_j$ from the decoder, therefore bound the difference between the actual output and $\alpha_s B_{c_1}\cdots B_{c_\ell}E D_j$ by
\[
                           D_*(\ell+1)\delta.
\]
Wordwise correctness of the original machine implies $e_{k,m}\le\varepsilon+\eta+D_*(m+1)\delta<\varepsilon+\gamma$, a contradiction. The representative $E$ is used only as part of the terminal decoder. No idempotence of $E$, consistency between different filler words, or hidden-state identity on a physical return was assumed.

The same argument shows that one may take
\begin{equation}\label{eq:ramseyloss56}
              L_{k,\varepsilon}=2g+(g+1)R_q(m+2).
\end{equation}
For exact cofinite returns, use exact fillers and omit $\eta$ and the neighborhood choice; $g$ is the given exact-return bound. In general $g$ depends on the finite witness and the continuity neighborhood. The remaining parameters depend only on the interface, the alphabet, $k$, and $\varepsilon$, never on a machine or horizon.

If $M_\varepsilon=\infty$, \eqref{eq:occupation56} for each $k$ gives $W_{N,\varepsilon}\to\infty$. If $M_\varepsilon=K<\infty$, the stationary machine gives $W_{N,\varepsilon}\le K$ for every $N$, and the bound for $k=K-1$ gives eventual equality when $K>1$. When $K=1$ equality holds for every $N$ because registers are nonempty.
\end{proof}

\begin{proposition}[Approximate returns without exact returns]\label{prop:approxreturns56}
Let the two circle commands have angles $\theta,\phi$, where $1,\theta/(2\pi),\phi/(2\pi)$ are rationally independent. They have no nonempty exact-return word, but satisfy Definition~\ref{def:returns56}.
\end{proposition}
\begin{proof}
A length-$n$ product has angle $n\phi+j(\theta-\phi)$, $0\le j\le n$. The positive orbit of the irrational rotation $\theta-\phi$ is dense. Compactness supplies, for every accuracy, a finite initial orbit segment that is a net of the circle. All longer initial segments remain nets; rotation by $n\phi$ preserves their accuracy. Thus a product is arbitrarily close to the identity at every sufficiently large length. An exact return would be a nontrivial integer relation among $1,\theta/(2\pi),\phi/(2\pi)$.
\end{proof}
A single irrational letter fails the eventual condition: there is only one word at each horizon. Its target can be placed in a one-label horizon-dependent decoder even when no finite stationary approximation at the same tolerance exists. Thus mere recurrence along a subsequence is not a substitute for Definition~\ref{def:returns56}.

\begin{remark}[Quantitative interpretation]
The Ramsey estimate gives a finite occupation constant, not a useful
asymptotic width rate. The homogeneous neutral representative remains
inside the terminal decoder in the extraction proof. It is neither an
identity transition nor an assumed idempotent.
\end{remark}

\section{Stationary stochastic purification}\label{sec:purification55}
Write $\mathsf{Stoch}(k)$ for the compact semigroup of row-stochastic matrices of order $k$. The identity of a group inside this semigroup can be a singular idempotent; it must not be confused with the identity matrix.

\begin{lemma}[An idempotent over the physical identity]\label{lem:idempotent55}
Let $\mathcal T$ be a compact subsemigroup of $H\times\mathsf{Stoch}(k)$ whose projection onto $H$ is surjective. There is an element $e=(1,E)\in\mathcal T$ with $E^2=E$ and rank equal to the minimum matrix rank in $\mathcal T$. The monoid $e\mathcal T e$ is a compact group with identity $e$ and still projects onto $H$.
\end{lemma}
\begin{proof}
Only the integer ranks $1,\ldots,k$ occur, so their nonempty set has a minimum which is attained. Choose $(h,T)$ of that minimum rank $r$. The powers of $T$ are bounded, since they are stochastic. Thus its complex eigenvalues have modulus at most one and all Jordan blocks on the unit circle have size one: a larger peripheral Jordan block would produce polynomially growing powers. The closure of positive powers of an element of a compact group is a group: arbitrarily large positive powers approach its identity, and the preceding powers approach its inverse. Apply this observation to the compact monothetic group generated by $(h,\lambda_1,\ldots,\lambda_b)$ in $H\times(S^1)^b$, where the $\lambda_i$ are the peripheral eigenvalues of $T$. There is a net of integers tending to infinity for which $h^n\to1$ and all peripheral eigenvalues raised to the $n$th power tend to one. Along it $T^n$ tends to the spectral projection $E$ onto the peripheral subspace. This is a real stochastic idempotent. Its rank is at most $r$ and at least $r$ by minimality, so it is $r$. Nets avoid any metrizability assumption on $H$.

Every matrix $B$ in $e\mathcal T e$ satisfies $B=EBE$ and has rank $r$. Its action on the row space of $E$ is consequently invertible. Restriction embeds $e\mathcal T e$ as a compact submonoid of $H\times\mathrm{GL}(r,\R)$, with identity $(1,I_r)$; injectivity follows because the restriction determines $(xE)B=xB$ for every row $x$, by $B=EBE$, and hence determines every row of $B$. A compact submonoid of a topological group is a subgroup, by the same positive-power return argument. The inverse images of its inverses are still stochastic. Finally, sandwiching an element with physical coordinate $h$ between $e$'s leaves that coordinate unchanged. Surjectivity follows.
\end{proof}

\begin{lemma}[The recurrent simplex]\label{lem:simplex55}
If $E$ is a stochastic idempotent of rank $r$, then
\[
 \Delta_k E=\{p\in\Delta_k:pE=p\}
\]
is a simplex with exactly $r$ vertices. Every group of stochastic matrices with identity $E$ permutes these vertices.
\end{lemma}
\begin{proof}
Regard $E$ as a transition matrix. Its closed communicating classes each have a unique stationary probability row supported on that class. Every stationary row is a convex combination of these rows: mass on a transient class is zero, and stationarity within a closed irreducible class determines a one-dimensional space. For completeness, transient mass vanishes because each transient state has a positive probability of reaching a closed class in a bounded number of steps; a common bound over the finitely many states makes transient mass decrease by a fixed factor unless it is zero. The closed-class rows have disjoint supports, hence are linearly independent.

Each row of $E$ is stationary because $E^2=E$. Conversely, every stationary row equals its product with $E$. The span of the closed-class stationary rows is therefore precisely the row space of $E$, so there are $r$ classes and $r$ simplex vertices. If $B$ has a stochastic inverse relative to $E$, multiplication by $B$ and by this inverse are inverse affine maps of this simplex. Affine bijections permute its vertices.
\end{proof}

\begin{theorem}[Purification without a width or error loss]\label{thm:purification55}
For every compact-group interface in Definition~\ref{def:stationary55}, a stationary $k$-label realization of error at most $\varepsilon$ admits a permutation realization of error at most $\varepsilon$ on at most $k$ labels. There is no return-word hypothesis and no restriction on the number of connected components of $H$.
\end{theorem}
\begin{proof}
For the given machine form
\[
 \mathcal T=\overline{\{(u_w^{-1},T_w):w\in\mathcal A^*\}}
       \subset H\times\mathsf{Stoch}(k).
\]
The inverse on the physical coordinate is important: concatenation satisfies $u_{vw}^{-1}=u_v^{-1}u_w^{-1}$ and $T_{vw}=T_vT_w$. Thus $\mathcal T$ is a compact semigroup. Positive-word density makes its projection onto $H$ surjective. By continuity, its every element $(h,B)$ satisfies
\begin{equation}\label{eq:closederror55}
             \norm{\alpha_sBD_j-F_{s,j}(h^{-1})}_j\le\varepsilon
             \quad(s\in\mathcal S,\ j\in\mathcal J).
\end{equation}

Apply Lemma~\ref{lem:idempotent55} and denote the vertices of $\Delta_k E$ by $\nu_1,\ldots,\nu_r$. For each letter the element
\[
        e(u_a^{-1},T_a)e=(u_a^{-1},ET_aE)
\]
belongs to the group $e\mathcal T e$. By Lemma~\ref{lem:simplex55}, $B_a=ET_aE$ permutes the $\nu_i$. Use that permutation as the new command. If $V$ has rows $\nu_i$ and the induced row permutation is $\Pi_a$, the precise intertwining identity is $VB_a=\Pi_aV$. Write uniquely
\[
          \alpha_sE=\sum_{i=1}^r\beta_s(i)\nu_i,
          \qquad \widetilde D_j(i)=\nu_iD_j\in Y_j.
\]
The last inclusion uses convexity of the legal output set. For a nonempty word the output of the resulting $r$-label machine is $\alpha_s B_{a_1}\cdots B_{a_n}D_j$. For the empty word it is $\alpha_sED_j$. In both cases the corresponding joint element lies in $\mathcal T$ and has physical coordinate $u_w^{-1}$. Equation~\eqref{eq:closederror55} proves the error bound. Since $r\le k$, the width has not increased.
\end{proof}
The new machine need not reproduce the old machine's erroneous output word by word; it remains in the same closed error balls about the target. Nor do we assert $T_{w}=I$ on a physical return, invertibility of the original $T_a$, or deterministic initialization. These three incorrect shortcuts would invalidate the reduction.

\subsection{Finite quotients and the unrestricted stationary minimum}
For a tuple of permutation matrices $\sigma=(\Pi_a)_{a\in\mathcal A}$ of order $k$, set
\[
 \mathcal L_\sigma=
 \overline{\{(u_w^{-1},\Pi_w):w\in\mathcal A^*\}}
       \subset H\times\mathfrak S_k.
\]
The topology is the product of the compact topology on $H$ and the discrete topology on the finite group $\mathfrak S_k$. This is a compact subgroup. Its definition includes the possible discrepancy between physical relations and hidden-state relations.

\begin{theorem}[An exact stationary minimum]\label{thm:stationaryminimum55}
The value $M_\varepsilon$ is the least positive integer $k$ for which some tuple $\sigma$ of permutations of order $k$, some $\beta_s\in\Delta_k$, and some $d_j(i)\in Y_j$ satisfy
\begin{equation}\label{eq:permutationfeasibility55}
 \norm{\beta_s\Pi d_j-F_{s,j}(h^{-1})}_j\le\varepsilon
 \quad\text{for all }(h,\Pi)\in\mathcal L_\sigma\text{ and all }s,j.
\end{equation}
For fixed $k$ there are at most $(k!)^{|\mathcal A|}$ discrete choices (a crude bound before simultaneous conjugacy reduction), and the remaining feasible set is compact. The optimal stationary error at any fixed width is attained.
\end{theorem}
\begin{proof}
Necessity follows from Theorem~\ref{thm:purification55}, padding by unused labels if its output has fewer than $k$ states, and then taking the closure of word products. Sufficiency is the restriction of \eqref{eq:permutationfeasibility55} to words. The probabilities and legal decoders lie in a finite product of compact sets. Each displayed inequality defines a closed subset of that product. The worst error is continuous in these parameters, uniformly in the compact group variable, so its minimum exists. There are finitely many permutation tuples.
\end{proof}

\begin{theorem}[Finite observable quotient criterion]\label{thm:finitequotient55}
For arbitrary compact $H$ and the stated continuous interface,
\[
 M_\varepsilon<\infty
 \quad\Longleftrightarrow\quad
 r(U)\le\varepsilon\text{ for some open normal subgroup }U\le H.
\]
If such $U$ exists, $|\mathcal S|[H:U]$ deterministic-initialized permutation labels suffice. Conversely, a stationary $k$-label realization supplies such a $U$ with $[H:U]\le k!$.
\end{theorem}
\begin{proof}
Purify a stationary realization and use its joint group $\mathcal L_\sigma$. Put
\[
                  U=\{h\in H:(h,I)\in\mathcal L_\sigma\}.
\]
This is a closed normal subgroup of $H$: conjugate $(h,I)$ by an element above any $g\in H$, using surjectivity of the first projection. For a fixed permutation $\Pi$, a nonempty fiber in $\mathcal L_\sigma$ is a coset $hU$. These at most $k!$ fibers cover $H$, so $U$ has finite index. A closed subgroup of finite index in a compact group is open, since its complement is a finite union of closed cosets.

For each $s,j$ and each fiber, the single legal vector $\beta_s\Pi d_j$ is within $\varepsilon$ of every $F_{s,j}(h^{-1})$ in that fiber. Explicitly, $(hU)^{-1}=Uh^{-1}=h^{-1}U$ by normality, so inversion preserves the ordinary $U$-coset form. Hence $r(U)\le\varepsilon$. Conversely, store $(s,gU)$, update the coset by left multiplication by $u_a$, and decode a minimizing center of $F_{s,j}(gU)$. Compactness supplies the centers. This gives the asserted finite realization.
\end{proof}
For finitely many connected components, every open subgroup contains $K=H^\circ$, and $K$ itself is open normal. Therefore $M_\varepsilon<\infty$ holds exactly when
\[
 \varepsilon\ge\max_{s,j,c\in H/K}\rad_{Y_j}(F_{s,j}(c)),
\]
\emph{without} a return assumption. Executable returns enter the stronger clocked converse, not this stationary statement. The index bound $k!$ is an existence bound on a deterministic quotient, not a claim that the stationary minimum equals its index.

\begin{proposition}[Random initialization strictly saves states]\label{prop:C655}
Let $H=\Z/6\Z$, with an identity letter and a generator, one seed, and binary success probability
\[
 f(n)=\tfrac12+\tfrac14(-1)^n+\tfrac14\cos(2\pi n/3).
\]
At zero error the least deterministic component-quotient width is six, whereas the unrestricted stationary stochastic minimum is five.
\end{proposition}
\begin{proof}
The values at $n=0,\ldots,5$ are
\[
                  (1,\tfrac18,\tfrac58,\tfrac12,\tfrac58,\tfrac18).
\]
This sequence has least period six. A deterministic quotient of the cyclic component action preserving these outputs must therefore have six states. For a five-state realization take the disjoint union of a two-cycle and a three-cycle, both updating the index by $i\mapsto i+1$ modulo their respective lengths, initialize with mass $1/2$ at the zeroth vertex of each, and use decoder rows $(1,0)$ on the two-cycle and $(1,1/4,1/4)$ on the three-cycle. The identity command fixes all five labels. Direct averaging gives exactly $f(n)$, with legal binary probabilities.

If a stationary realization on at most four labels existed, purification would give one whose generator is a permutation of at most four letters. Its order is one of $1,2,3,4$, as follows by listing the possible cycle lengths. Every resulting numerical sequence has a period dividing that order, contradicting the least period six. Thus five is the unrestricted minimum.
\end{proof}

\section{Physical equivariance and the complete boundary}\label{sec:boundary56}
Return now to a compact Hausdorff group $H$ with a finite dense generating alphabet. The output sets may still be norm-compact convex subsets of real Banach spaces. Purification removes transition randomness; a further averaging removes hidden relations that have no physical counterpart.

\begin{theorem}[Physical-equivariant purification]\label{thm:equivariant56}
Every stationary $k$-label realization of error at most $\varepsilon$ admits one on $r\le k$ labels of the form
\begin{equation}\label{eq:equivariant56}
 \alpha_s\rho(h^{-1})d_j,\qquad h\in H,
 \qquad \rho:H\longrightarrow S_r\text{ a continuous homomorphism},
\end{equation}
with the same uniform error. Here permutations are represented by row-action matrices. In particular the command row is $\rho(u_a^{-1})$, all physical group relations hold on the new register, and initialization may remain random.
\end{theorem}
\begin{proof}
Apply Theorem~\ref{thm:purification55}, whose proof uses only finite convex combinations of decoder values. Denote the resulting initialization by $\beta_s$, its decoder by $D_j$, and its compact joint group by $\mathcal L\subset H\times S_k$, with the product topology and discrete permutation factor. Unused labels may be adjoined when necessary. Let $P$ be the second projection and
\[
                 N=\{\pi\in P:(1,\pi)\in\mathcal L\}.
\]
Then $N$ is normal in $P$. The averaging matrix $Q=|N|^{-1}\sum_{n\in N}n$ is stochastic, satisfies $Q^2=Q$, and commutes with $P$. For $(h,\pi)\in\mathcal L$, every $(h,n\pi)$ belongs to $\mathcal L$. Averaging their error inequalities and using convexity of the norm ball gives
\begin{equation}\label{eq:kernelaverage56}
                   \|\beta_sQ\pi D_j-F_{s,j}(h^{-1})\|_j\le\varepsilon.
\end{equation}

Let $O_1,\ldots,O_r$ be the $N$-orbits in the label set, and let $v_i$ be the uniform probability row on $O_i$. The rows of $Q$ are precisely these uniform orbit rows. Write $V$ for the matrix with rows $v_i$. Every $\pi\in P$ maps $N$-orbits bijectively onto $N$-orbits and therefore induces a permutation $\bar\pi$ with
\begin{equation}\label{eq:orbitintertwining56}
                             V\pi=\bar\pi V.
\end{equation}
The induced action of each $n\in N$ is trivial. Two elements of $\mathcal L$ with the same first coordinate differ by an element of $\{1\}\times N$. Hence $(h,\pi)\mapsto\bar\pi$ descends to a homomorphism $\rho:H\to S_r$. It is continuous: the surjection from the compact group $\mathcal L$ onto the Hausdorff group $H$ is a quotient map, and the original orbit action is continuous into a finite discrete group.

Set $\alpha_s(i)=\sum_{x\in O_i}\beta_s(x)$ and $d_j(i)=v_iD_j$. Then $\alpha_sV=\beta_sQ$ and $d_j=VD_j$. Equation~\eqref{eq:orbitintertwining56} gives
\[
                 \alpha_s\rho(h)d_j=\beta_sQ\pi D_j
                 \quad((h,\pi)\in\mathcal L).
\]
Replace $h$ by $h^{-1}$ in \eqref{eq:kernelaverage56}. This proves \eqref{eq:equivariant56}; all decoder values are legal by convexity, and $r\le k$.
\end{proof}

\begin{corollary}[Intrinsic stationary minimum]\label{cor:intrinsicminimum56}
The number $M_\varepsilon$ is the least cardinality of a finite continuous $H$-set for which some probability initialization for each seed and legal decoder for each query approximate the interface within $\varepsilon$. It is independent of the chosen finite dense alphabet. It is unchanged by replacing every target with
\[
                       h\longmapsto F_{s,j}(\varphi(h)q),
\]
where $\varphi$ is a continuous group automorphism and $q\in H$.
\end{corollary}
\begin{proof}
Theorem~\ref{thm:equivariant56} proves necessity, and restriction to command words proves sufficiency. For the last assertion, replace $\rho$ by $\rho\circ\varphi$ and $\alpha_s$ by $\alpha_s\rho(q^{-1})$ in \eqref{eq:equivariant56}. Applying the inverse change of variables proves equality of the minima.
\end{proof}
The finite quotient in Theorem~\ref{thm:finitequotient55} has a more precise description than its factorial bound. With $\mathcal L,P,N$ as above and $U=\{h:(h,I)\in\mathcal L\}$, the correspondence $(h,\pi)\mapsto(hU,\pi N)$ induces $H/U\simeq P/N$. Thus $[H:U]=|P/N|\le k!$. The orbit action in Theorem~\ref{thm:equivariant56} can have an additional kernel; no equality between this index and the minimal label count is asserted.

\begin{theorem}[Complete clocked finite-quotient classification]\label{thm:completeboundary56}
Assume eventual approximate returns on the compact group $H$. At every $\varepsilon\ge0$, including a positive critical radius, the following are equivalent:
\begin{enumerate}
\item $\sup_NW_{N,\varepsilon}<\infty$;
\item $M_\varepsilon<\infty$;
\item $r(U)\le\varepsilon$ for an open normal subgroup $U$ of $H$.
\end{enumerate}
For each $k<M_\varepsilon$ the all-profile occupation bound \eqref{eq:occupation56} holds, and $W_{N,\varepsilon}$ converges to $M_\varepsilon$. Put
\[
 R_0=\max_{s,j,h\in H}\rad_{Y_j}(F_{s,j}(hH^\circ)).
\]
Then $R_0=\inf_{U\lhd H\text{ open}}r(U)$. Above $R_0$ bounded realizations exist; below it every fixed width has bounded occupation. At $\varepsilon=R_0$ bounded clocked realization exists exactly when the infimum is attained by an open normal subgroup.
\end{theorem}
\begin{proof}
The first two conditions are equivalent by Theorem~\ref{thm:stationarization56}. The second and third are equivalent by Theorem~\ref{thm:finitequotient55}. That proof, including constrained center attainment, uses compactness and convexity of $Y_j$, not finite dimensionality.

For completeness, every open subgroup contains $H^\circ$: the image of the connected identity component in a finite discrete quotient is a singleton. Thus $r(U)\ge R_0$. The quotient $H/H^\circ$ is profinite, so its open normal subgroups form a neighborhood basis at the identity. The inverse images $U$ have intersection $H^\circ$. Given a neighborhood $O$ of $H^\circ$ in $H$, compactness supplies one such $U\subset O$: otherwise their intersections with the closed set $H\setminus O$ have the finite intersection property. Uniform continuity of the finite family of targets gives, for every $\eta>0$, a neighborhood $V$ of $1$ such that $\|F_{s,j}(hkv)-F_{s,j}(hk)\|_j<\eta$ for all $h\in H$, $k\in H^\circ$, and $v\in V$. Choose $U\subset H^\circ V$. Every $F_{s,j}(hU)$ is then within $\eta$ of $F_{s,j}(hH^\circ)$, so $r(U)\le R_0+\eta$. This proves the infimum formula and all remaining assertions.
\end{proof}
This result includes the formerly separate zero and positive boundary cases. The finite-rank proof of Theorem~\ref{thm:zero55} is retained below because it identifies an additional exact translate-space obstruction. No passage from an error strictly larger than $R_0$ to equality is used in Theorem~\ref{thm:completeboundary56}.

\subsection{Length phases without a width overhead}
Let $p\ge1$ be an integer and fix a command with product $a$. Let $J$ be the closure of all products of lengths divisible by $p$. The group argument of Lemma~\ref{lem:phasegroup55}, which uses only the block length, shows that $J$ is open normal, $H/J$ is cyclic of order dividing $p$, and all letters belong to $Ja$. Regard $\mathcal A^p$ as a finite alphabet on $J$. Assume that this block alphabet has eventual approximate returns. In particular this holds when the original nonempty exact-return language has greatest common divisor $p$.

For $0\le r<p$ let $M_{r,\varepsilon}$ be the intrinsic stationary minimum on $J$ for
\[
                         f_{r,s,j}(h)=F_{s,j}(ha^r).
\]
\begin{theorem}[Exact limiting width in each phase]\label{thm:phaseminimum56}
Under the preceding assumptions, for every $r$ and $\varepsilon\ge0$,
\begin{equation}\label{eq:phaseminimum56}
 W_{N,\varepsilon}\le M_{r,\varepsilon}\quad(N\equiv r\pmod p),
 \qquad
 \lim_{\substack{N\to\infty\\N\equiv r\ ({\rm mod}\ p)}}W_{N,\varepsilon}
          =M_{r,\varepsilon}.
\end{equation}
Finite limits are eventually attained. If $k<M_{r,\varepsilon}$, the number of positive cuts of width at most $k$ is bounded uniformly over that horizon phase, with no restriction on other widths. Boundedness in the phase is equivalent to the existence of an open normal subgroup $U\le J$ such that
\begin{equation}\label{eq:phasequotient56}
          \max_{s,j,h\in J}\rad_{Y_j}(F_{s,j}(hUa^r))\le\varepsilon.
\end{equation}
\end{theorem}
\begin{proof}
Suppose first that $M_{r,\varepsilon}=K<\infty$ and choose a physical-equivariant stationary realization of $f_r$ on $K$ labels, with representation $\rho:J\to S_K$. For a fixed $N\equiv r\pmod p$, let $g_t$ be the physical product after $t$ commands and set
\begin{equation}\label{eq:gauge56}
 z_t=a^{N-t}g_ta^{-r}\in J,
 \qquad c_{t,b}=a^{N-t}u_ba^{-(N-t+1)}\in J.
\end{equation}
Normality of $J$ and $u_ba^{-1}\in J$ justify the membership assertions. We have $z_0=a^{N-r}$, $z_t=c_{t,b}z_{t-1}$ when the $t$th letter is $b$, and $z_N=g_Na^{-r}$. Initialize with $\alpha_s\rho(z_0^{-1})$, use the clock-dependent permutation $\rho(c_{t,b}^{-1})$, and keep decoder $d_j$. The row product is $\alpha_s\rho(z_N^{-1})$, so its error from $F_{s,j}(g_N)=f_{r,s,j}(z_N)$ is at most $\varepsilon$. This uses the free external clock but no phase register or partial-block storage. It proves the upper bound with exactly $K$ labels.

Fix $k<M_{r,\varepsilon}$ and a cut residue $0\le\ell<p$. Put $b=(r-\ell)\bmod p$, chosen in $\{0,\ldots,p-1\}$. Restrict the input to the fixed prefix $a^\ell$, arbitrary $p$-letter blocks, and fixed suffix $a^b$. Integrate the prefix into initialization and the suffix into the decoder. The block target on $J$ is
\[
 G_{\ell,s,j}(h)=F_{s,j}(a^b h a^\ell)
     =f_{r,s,j}\bigl(\operatorname{Ad}_{a^b}(h)a^{\ell+b-r}\bigr).
\]
Here $a^{\ell+b-r}\in J$ and conjugation by $a^b$ is an automorphism of $J$. Corollary~\ref{cor:intrinsicminimum56} therefore identifies its stationary minimum with $M_{r,\varepsilon}$. Theorem~\ref{thm:stationarization56} on the block alphabet bounds its positive narrow cuts by a constant $L_\ell$. These cuts are precisely the original cuts in residue $\ell$ between the prefix and suffix, apart from the initial block cut. At most  $\ell+b+1\le2p$ further cuts need be charged for this residue, including small horizons shorter than the chosen prefix and suffix. Summing gives the uniform bound $\sum_{\ell=0}^{p-1}L_\ell+2p^2$. As before, a peak at most $k$ would force $N$ below this bound. This proves the limit and occupation assertions. Finally apply Theorem~\ref{thm:finitequotient55} to $f_r$ on $J$ to obtain \eqref{eq:phasequotient56}.
\end{proof}
If the order of $H/J$ is a proper divisor $q$ of $p$, the minima and quotient-radius criteria repeat with period $q$: multiplication by $a^q\in J$ is a right translation of the target. The earlier phase radii have the same repetition. Neither the physical quotient order nor the return-length greatest common divisor is substituted silently for the other.

\subsection{Compact outputs and finite hidden state}
Theorems~\ref{thm:stationarization56}, \ref{thm:equivariant56}, \ref{thm:completeboundary56}, and \ref{thm:phaseminimum56} hold for norm-compact convex Banach-valued outputs. A single-machine purification even needs only convexity: its new decoder values lie in the convex hull of finitely many old values. Compact output sets are used instead in finite-witness extraction, optimal-error attainment, and constrained-center existence. The representative-function and algebraic sections below retain their explicitly finite-dimensional hypotheses.

Finiteness on the hidden side has a different role. On countably many labels the full simplex is not compact in total variation; mass can escape to infinity, matrix ranks need not admit a useful finite minimum, and a recurrent simplex may have infinitely many vertices. On a general compact convex hidden state space affine automorphisms need not permute a finite vertex set: rotation of the disk already acts continuously on infinitely many extreme points. Thus these proofs give neither a finite-label purification theorem for arbitrary operator kernels nor an online observation theorem. Their terminal-output assertion is strictly the one stated in the models.

\begin{remark}[Action and counting conventions]
Probability vectors act on the right of transition matrices. Chronological
row multiplication consequently corresponds to the inverse physical
coordinate $h^{-1}$ used above. The coarse deterministic quotient
construction stores both the seed and the coset; its
$|\mathcal S|[H:U]$ bound is not the unrestricted stochastic minimum.
In the phase theorem the chosen block period $p$, the physical quotient
order $q\mid p$, and the greatest common divisor of actual return
lengths retain their separate meanings.
\end{remark}

\section{The clocked component-radius threshold}\label{sec:intro54}
A finite stochastic register can carry infinitely many probability distributions. Thus an infinite numerical orbit is not, by itself, an obstruction to finite positive realization. A further difficulty arises when the realization may be redesigned at every horizon and its transitions may depend on a free external clock. The obstruction established here survives both freedoms. It is measured by the radius of an observable image inside its allowed output set.

For the clocked theorem in this section, let $H$ be a compact Hausdorff group with finitely many connected components. Put $K=H^\circ$ and $C=H/K$. Then $K$ is an open normal subgroup and $C$ is a finite group. Fix a finite alphabet $\mathcal A$ with elements $u_a\in H$ whose generated group is dense in $H$. Our convention is
\[
 u_{a_1\cdots a_n}=u_{a_n}\cdots u_{a_1}.
\]
The identity $e$ need not be an alphabet element. Instead assume the following executable return condition:
\begin{equation}\label{eq:returns54}
 \text{there is an integer }g\ge0\text{ such that, for every }n\ge g,
 \text{ some word of length }n\text{ has product }e.
\end{equation}
An identity letter gives $g=0$. The condition concerns physical products; the stochastic rows on a return word may perform arbitrary processing.

For nonempty finite seed and query sets $\mathcal S,\mathcal J$, let
\[
 F_{s,j}:H\longrightarrow Y_j
\]
be continuous, where $Y_j$ is a nonempty compact convex subset of a finite-dimensional real normed space $(V_j,\norm{\cdot}_j)$. A query is selected only after the command word. A prescribed joint law of finitely many outputs is treated as one query, not as a collection of marginal constraints. In particular, a categorical interface has
\[
 Y_j=\Delta_{M_j}=\{p\in\R^{M_j}:p_i\ge0,\ \textstyle\sum_i p_i=1\},
 \qquad \norm{x}_j=\tfrac12\sum_i|x_i|.
\]
Then $\norm{p-q}_j$ is total variation. Binary means correspond to $Y=[-1,1]$ with $\norm{x}=|x|/2$.

\begin{definition}[Clocked realization]\label{def:machine54}
At horizon $N$ choose finite nonempty registers $S_t$, $0\le t\le N$, initialization rows $E_s$, row-stochastic command matrices $T_{t,a}:S_{t-1}\to S_t$, and terminal decoder maps $D_j:S_N\to Y_j$. The realized output vector on $w=a_1\cdots a_N$ is
\[
 E_sT_{1,a_1}\cdots T_{N,a_N}D_j.
\]
Its error is the maximum, over all seeds, length-$N$ words, and queries, of its distance from $F_{s,j}(u_w)$. The command width is $\max_t|S_t|$, counting all declared labels, including unused padding labels. Write $W_{N,\varepsilon}$ for its least possible value at error at most $\varepsilon$.
\end{definition}
All retained private information, including randomness kept for future use, is charged to the current label. The rows may depend on $t$ and $N$, but not on an unrecorded input history or on the future query. The external clock and horizon, arbitrary real table descriptions, their construction and lookup, real arithmetic, and exact sampling are uncharged. This is a numerical realization invariant, not uniform workspace or program size. For a finite categorical interface, charging the terminal answer register changes the peak to $\max(W_{N,\varepsilon},\max_j M_j)$. Thus all boundedness and growth results are unchanged; the convention that also charges a sampled binary-answer register has a minimum of two labels. A decoder vector by itself is not an additional sampled register.

For $A\subset Y_j$ define its constrained radius by
\[
 \rad_{Y_j}(A)=\min_{y\in Y_j}\sup_{x\in A}\norm{x-y}_j.
\]
The minimum is attained when $A$ is compact. The critical error is
\begin{equation}\label{eq:radius54}
 \epsc=\max_{s,j,c\in C}\rad_{Y_j}\bigl(F_{s,j}(c)\bigr).
\end{equation}
Here a component is identified with its coset in $H$; its image is compact by continuity. The requirement that the center belong to $Y_j$ is essential for categorical outputs.

\begin{theorem}[Intrinsic radius classification]\label{thm:radius54}
Under \eqref{eq:returns54}, for every $\varepsilon\ge0$ the following are equivalent:
\begin{enumerate}
\item $\sup_N W_{N,\varepsilon}<\infty$;
\item one stationary finite-state realization with permutation command updates has error at most $\varepsilon$ at every horizon;
\item $\varepsilon\ge\epsc$.
\end{enumerate}
At the boundary, at most $|\mathcal S||C|$ command labels suffice. If $\varepsilon<\epsc$, then for every integer $k\ge1$ there is an integer $L_k<\infty$ such that every error-$\varepsilon$ realization satisfies
\begin{equation}\label{eq:occupation54}
 \#\{t\in\{1,\ldots,N\}:|S_t|\le k\}\le L_k.
\end{equation}
All intervening widths are unrestricted. In particular, $W_{N,\varepsilon}\to\infty$.
\end{theorem}
Neither a positive lower bound on state masses nor a reachable-rank, observability, or conditioning assumption occurs in the theorem. The upper realization is deterministic until its terminal decoder. It need not have the minimum width among all stochastic realizations. Theorem~\ref{thm:stationaryminimum55}, rather than a deterministic component partition, characterizes that stationary minimum.

For binary means, \eqref{eq:radius54} reduces to one quarter of the largest componentwise oscillation. For a genuine categorical law it is generally not half its diameter; an analytic three-outcome example below has critical total-variation error $2/3$, whereas every individual event has binary threshold $1/2$. The proof uses finitely many localized vector tests rather than a separate midpoint for each coordinate.

The finite-component assumption does not assert that $H$ is a Lie group. For example, take $H=\prod_{n\ge1}\R/\Z$, choose an element whose every finite coordinate tuple together with $1$ is rationally independent, and include this element and $e$ as commands. The positive orbit is dense and $H$ is connected but has no faithful finite-dimensional representation. The theorem applies to every finite continuous interface on this group. With infinitely many components the component construction need not be finite; Section~\ref{sec:profinite55} gives the finite-observable-quotient criterion and the strict-side clocked theorem in that setting.

The quantitative orbit and arithmetic results have separate hypotheses and are treated in the quantitative paper and the complete research edition.

\subsection{Executable tests}
Let $\mathcal P=\{u_w:w\in\mathcal A^*\}$. We use standard Haar averaging and the density of matrix coefficients of finite-dimensional unitary representations, in realified form \cite{Folland}. No quantitative mixing assumption for a preassigned random walk is imposed.

\begin{lemma}[Positive words and common lengths]\label{lem:positive54}
The semigroup $\mathcal P$ is dense in $H$, and $\mathcal P\cap K$ is dense in $K$. Every finite subset of $\mathcal P\cap K$ has executable representatives of a common positive length.
\end{lemma}
\begin{proof}
For an element $u$ of a compact group, positive powers return arbitrarily close to $e$ with exponent at least two. Indeed cover the compact closure of its powers by finitely many translates of a sufficiently small neighborhood. Among sufficiently many powers, three lie in one member of the cover; the first and last exponents differ by at least two, and their quotient is arbitrarily close to $e$. This argument also covers finite-order elements by taking a sufficiently large multiple of the order. If $u^r$ approaches $e$ with $r\ge2$, then $u^{r-1}=u^ru^{-1}$ approaches $u^{-1}$ and is still a positive power. Thus $u^{-1}$ is in the closure of positive powers. The closed semigroup $\overline{\mathcal P}$ therefore contains the group generated by the letters, and hence equals $H$. Since $H$ has finitely many components, $K$ is open, and a relatively open subset of $K$ is open in $H$; density gives the second assertion. For finitely many representative lengths $\ell_i$, choose $B\ge\max_i\ell_i+g$, $B\ge1$, and append return words of lengths $B-\ell_i$. The products are unchanged.
\end{proof}

\begin{lemma}[Finite executable gap]\label{lem:gap54}
Let $R:K\to O(E)$ be a continuous representation on a nonzero finite-dimensional Euclidean space, with no nonzero fixed vector. For every $k\ge1$ there are $B\ge1$, a probability law $\nu$ on finitely many executable length-$B$ words with products in $K$, and $d\in(0,1)$ such that
\begin{equation}\label{eq:gap54}
 \E_{w\sim\nu}\max_{1\le i\le k}\ip{v_i}{R(u_w)u}\le1-d
\end{equation}
for all unit vectors $u,v_1,\ldots,v_k\in E$.
\end{lemma}
\begin{proof}
On the compact product of unit spheres consider the minimum
\[
 \gamma=\min_{u,v_1,\ldots,v_k}\int_K
 \left(1-\max_i\ip{v_i}{R(h)u}\right)\,d\mu(h).
\]
Repeated centers are permitted. The integrand is nonnegative. If a minimizing tuple had integral zero, full support of Haar measure would give $R(h)u\in\{v_1,\ldots,v_k\}$ for every $h$. The orbit of $u$ is connected, hence a singleton, contradicting the fixed-vector hypothesis. Thus $\gamma>0$.

Partition the compact image $R(K)$ into finitely many measurable sets of small diameter in the Euclidean operator norm, and replace their Haar masses by atoms at nearby executable products in $K$. Cells of zero Haar mass can be discarded. The function $A\mapsto\max_i\ip{v_i}{Au}$ is uniformly $1$-Lipschitz in that norm. An approximation radius less than $\gamma/2$ therefore preserves a positive defect uniformly in the unit vectors. Lemma~\ref{lem:positive54} gives a common length. Decreasing the defect to a number in $(0,1)$ proves the assertion.
\end{proof}

\begin{lemma}[Conditional-centroid contraction]\label{lem:centroid54}
Let $Z\in E$ be a bounded random physical vector and $S$ the current register. Apply an independent word $w\sim\nu$ from Lemma~\ref{lem:gap54}, and let the actual block kernel generate a terminal state $S'$ having at most $k$ labels. Then
\begin{equation}\label{eq:centroid54}
 \E\norm{\E[R(u_w)Z\mid S']}\le(1-d)\E\norm{\E[Z\mid S]}.
\end{equation}
A fixed intervening return word cannot increase the same conditional norm, regardless of its intermediate or terminal widths.
\end{lemma}
\begin{proof}
Put $z_s=\E[Z\mid S=s]$ and write $T_w(s,i)$ for the actual composite row of the block. Since this row depends on the past only through $S$,
\[
 \Pr(S'=i)\E[R(u_w)Z\mid S'=i]
 =\sum_s\Pr(S=s)\E_w T_w(s,i)R(u_w)z_s.
\]
Choose $v_i$ in the direction of each nonzero terminal centroid, arbitrarily otherwise. Pair, sum in $i$, and use stochasticity to bound the sum by
\[
 \sum_s\Pr(S=s)\E_w\max_i\ip{v_i}{R(u_w)z_s}
 \le(1-d)\sum_s\Pr(S=s)\norm{z_s}.
\]
Zero-probability states contribute zero. No restriction was placed on registers inside the block. On a fixed return word the physical vector is unchanged at its endpoints; the new state is a stochastic image of the old one, so conditional Jensen gives nonincrease. Independence is used only at complete block boundaries.
\end{proof}

\section{Representative functions and localized radii}\label{sec:local54}
A real representative function on $K$ is a matrix coefficient of a finite-dimensional continuous orthogonal representation, or a finite sum of such coefficients. Their algebra contains the constants, is closed under multiplication by tensor products, and is uniformly dense in $C(K,\R)$ by the Peter--Weyl approximation theorem \cite{Folland}. A finite collection of representative functions can be written on one space
\[
 \mathcal V=\bigoplus_{i=1}^a\End(E_i),\qquad
 \Phi(h)=(R_i(h))_{i=1}^a,
\]
where the scalar product is the sum of Frobenius products. Left multiplication defines an orthogonal representation $R$ on $\mathcal V$ and satisfies $\Phi(gh)=R(g)\Phi(h)$. Adjoin a trivial summand when necessary. Let
\[
 P=\int_KR(h)\,d\mu(h),\quad E=(I-P)\mathcal V,\quad
 Z(h)=(I-P)\Phi(h),\quad Q_0=\norm{Z(e)}.
\]
Haar averaging is the orthogonal projection onto fixed vectors. Consequently $Z(gh)=R(g)Z(h)$, the representation on $E$ has no fixed vectors, and $P\Phi(h)=P\Phi(e)$ for $h\in K$.

Choose and fix a coefficient vector $a_p$ for each scalar representative function $p$, and a linear coefficient map $A_G$ for each vector-valued representative function $G$. Their norms need not be intrinsic; fixing any such choices suffices. For an arbitrary joint distribution of $h\in K$ and a finite state $S$, put $Q=\E\norm{\E[Z(h)\mid S]}$ and $\bar p=\int p\,d\mu$. The fixed/moving decomposition gives
\begin{align}
 \E[p(h)\mid S]&=\bar p+\ip{a_p}{\E[Z(h)\mid S]},\label{eq:calibration54}\\
 |\E p-\bar p|&\le\norm{a_p}Q,\qquad
 \norm{\E G-\bar G}\le\norm{A_G}Q.\label{eq:vectorcal54}
\end{align}
Here the norm of $A_G$ is the operator norm from the Euclidean feature space to the indicated output space. If $D:S\to Y$ and $M_Y=\max_{y\in Y}\norm{y}$, then
\begin{equation}\label{eq:decodercal54}
 \norm{\E[p(h)D(S)]-\bar p\E D(S)}\le M_Y\norm{a_p}Q.
\end{equation}
Indeed condition on $S$ in \eqref{eq:calibration54}, multiply by $D(S)$, and apply the triangle inequality. These are moment estimates, not total-variation convergence of a hidden physical distribution.

\begin{lemma}[Finite localization of a radius]\label{lem:radiuslocal54}
Let $f:K\to Y$ be continuous, with $Y$ compact and convex in a finite-dimensional normed space. Suppose $\rad_Y(f(K))>\varepsilon$. There are a vector-valued representative function $G$, a number $\eta>0$, and nonnegative representative functions $p_1,\ldots,p_a$ of Haar mean one such that
\[
 \norm{G-f}_\infty<\eta,\qquad
 r_*:=\min_{y\in Y}\max_i\norm{\bar{p_iG}-y}>\delta:=\varepsilon+\eta.
\]
\end{lemma}
\begin{proof}
Choose $\varepsilon<t<\rad_Y(f(K))$. For each $y\in Y$ some $h\in K$ has $\norm{f(h)-y}>t$. The corresponding open subsets of $Y$ cover $Y$, so a finite subcover yields $h_1,\ldots,h_a$ with
\[
 \min_{y\in Y}\max_i\norm{f(h_i)-y}>t.
\]
The strict inequality follows also at a minimizing $y$, since the maximum is everywhere greater than $t$ and $Y$ is compact.

Fix a sufficiently small positive $\eta$. Approximate the coordinates of $f$ uniformly by representative functions to obtain $\norm{G-f}_\infty<\eta$. For each $i$, choose a nonzero nonnegative continuous bump supported in a neighborhood where $f$ differs from $f(h_i)$ by less than $\eta$. Such a bump exists by normality of the compact Hausdorff space; it has positive Haar integral by full support. Approximate its square root uniformly by a representative function, square the approximation, and normalize its Haar integral. The resulting nonnegative representative function $p_i$ has mean one, and its weighted $f$-mean can be made arbitrarily close to that of the bump. Thus we may ensure
\[
 \norm{\bar{p_iG}-f(h_i)}<3\eta.
\]
The constrained radius changes by at most the maximum perturbation of the finitely many points. Taking $4\eta<t-\varepsilon$ gives $r_*>\varepsilon+\eta$.
\end{proof}

\begin{proposition}[Residual moment from a vector interface]\label{prop:residual54}
Fix the data of Lemma~\ref{lem:radiuslocal54}. Choose one feature space representing all $p_i$ and $p_iG$, and fix their coefficient choices before defining
\[
 L=\max_i\bigl(\norm{A_{p_iG}}+(M_Y+\delta)\norm{a_{p_i}}\bigr),
 \qquad \Delta=r_*-\delta>0.
\]
Suppose a terminal decoder $D(S)\in Y$ obeys
\[
 \norm{\E[D(S)\mid w]-G(h(w))}\le\delta
\]
for every complete deterministic word in a test law. Then $L>0$, $E\ne\{0\}$, $Q_0>0$, and
\begin{equation}\label{eq:residual54}
                         Q\ge\Delta/L.
\end{equation}
\end{proposition}
\begin{proof}
For $p=p_i\ge0$, wordwise correctness and the triangle inequality imply
\[
 \norm{\E[p(h)D(S)]-\E[p(h)G(h)]}\le\delta\E p(h).
\]
This is a conditional mean error, not a pointwise bound on the random decoder. Since $\bar p_i=1$, calibration gives explicitly
\[
                  \E p_i\le1+\norm{a_{p_i}}Q.
\]
Write $y=\E D(S)$, which belongs to $Y$ by convexity. Equations~\eqref{eq:vectorcal54}--\eqref{eq:decodercal54}, together with $\bar p=1$, give
\[
 \norm{y-\bar{p_iG}}\le\delta+
 \bigl(\norm{A_{p_iG}}+(M_Y+\delta)\norm{a_{p_i}}\bigr)Q.
\]
Maximizing in $i$ yields $r_*\le\delta+LQ$. Hence $LQ\ge\Delta>0$. In particular the moving feature and its initial norm cannot vanish: by equivariance, $Q_0=0$ would imply $Z(h)=0$ for all $h$ and therefore $Q=0$.
\end{proof}

\subsection{Proof of Theorem~\ref{thm:radius54}}
Choose $(s,j,c)$ with component radius greater than $\varepsilon$. By density there is an executable prefix $v$ with $u_v\in c$. Put $b=|v|$ and $f(h)=F_{s,j}(hu_v)$ on $K$. Since $Ku_v=c$, this map has the same constrained radius. Apply Lemma~\ref{lem:radiuslocal54} and Proposition~\ref{prop:residual54}.

The moving feature is nonzero independently of any particular test law. Otherwise all represented functions would be constant on $K$; since each $p_i$ has mean one, this would force $p_i=1$ and $G$ constant. Then all $\bar{p_iG}$ would equal a point at distance less than $\eta$ from $Y$, contradicting $r_*>\delta\ge\eta$.

For a fixed width $k$, apply Lemma~\ref{lem:gap54} to the resulting moving feature. Let $B,d$ be its block parameters and set
\[
 \chi=-\log(1-d)>0,\qquad A=\max(1,LQ_0/\Delta).
\]
Start a test word with $v$; between independent length-$B$ test blocks use fixed return words. The accumulated physical product is $hu_v$. Initially $h=e$, so the conditional feature norm after the prefix equals $Q_0$, regardless of the hidden-state distribution. If $J$ block endpoints have width at most $k$, Lemma~\ref{lem:centroid54} gives
\[
 Q_N\le Q_0(1-d)^J.
\]
The machine is correct on every deterministic word of length $N$. Thus any external law supported on such words inherits the conditional mean bound, with error $\delta=\varepsilon+\eta$ for $G$. Proposition~\ref{prop:residual54} implies
\begin{equation}\label{eq:J54}
                         J\chi\le\log A.
\end{equation}
The law depends on the advertised width profile, not on the machine's private state. When $A=1$, \eqref{eq:J54} excludes even one selected narrow block; no logarithmic loss has occurred.

For clarity we give the spacing count with the full return hypothesis. Discard narrow cuts before $b+B+g$ and after $N-g$. From the remaining cuts greedily choose the first and then the next at distance at least $B+g$. Put an independent block immediately before each selected cut. The initial gap, all intervening gaps, and the final suffix have lengths at least $g$, so \eqref{eq:returns54} fills them with executable returns. At most $b+B+2g-1$ cuts were discarded. Each chosen endpoint accounts for at most $B+g$ eligible cuts. Therefore one may take
\begin{equation}\label{eq:Lk54}
 L_k=\left\lceil b+B+2g-1+(B+g)\frac{\log A}{\chi}\right\rceil.
\end{equation}
The constants $B,d$ depend on the fixed moving representation and $k$; $\chi=-\log(1-d)$ depends on that defect. The localization, prefix length $b$, and $A$ depend only on the alphabet, interface, chosen component, and error. Consequently $L_k$ depends on these fixed data, the return certificate $g$, and $k$, but not on $N$ or the machine. When the eligible interval is empty the same estimate follows directly by counting. This proves occupation without restricting any register inside a block or gap. If the entire peak is at most $k$, then $N\le L_k$, which proves divergence directly. Monotonicity in the horizon is not needed when an identity letter is absent.

For the upper bound, use labels $(s,c)\in\mathcal S\times C$, including any labels not reached at a given horizon. Initialize at $(s,K)$ and update by left multiplication
\[
 (s,c)\longmapsto(s,[u_a]c).
\]
Choose a minimizing center of $F_{s,j}(c)$ in $Y_j$ as decoder at $(s,c)$ for query $j$. These are legal decoder values, and the error is at most \eqref{eq:radius54}. All command updates are permutations of one finite state set, independent of clock and horizon. This proves boundary attainment and the equivalence.\qed

\begin{proposition}[A finite return certificate]\label{prop:return54}
Condition~\eqref{eq:returns54} follows if there are finitely many executable identity-product words whose positive lengths have greatest common divisor one.
\end{proposition}
\begin{proof}
Let the lengths be $\ell_1,\ldots,\ell_a$. Their nonnegative sums modulo $\ell_1$ form the subgroup generated by their residues: a subsemigroup of a finite group is a subgroup. The gcd hypothesis makes this the whole residue group. Choose one nonnegative sum $n_r$ for each residue $r$ modulo $\ell_1$, and put $g=\max_r n_r$. Every $n\ge g$ equals its representative $n_r$ plus a nonnegative multiple of $\ell_1$. Concatenating the corresponding return words proves the claim.
\end{proof}
For example, the alphabet $\{R_\theta,R_{-\theta},R_{2\pi/3}\}$, with $\theta/2\pi$ irrational, contains no identity letter but has return lengths two and three. It satisfies \eqref{eq:returns54} with $g=2$. A single irrational rotation letter does not: there is only one word at each horizon, and a horizon-specific constant decoder can store its numerical answer with one command label. Thus synchronization cannot be silently omitted. Nor is a bounded irreversible semigroup enough: for the identity command $I=1$ and contraction command $a=\lambda\in(0,1)$, the mean $\rho\lambda^{\#a(w)}$ has an exact alive/absorbed two-state realization. The group generated with inverses is not bounded. The present theorem is not a compact-semigroup classification.

\section{Exact-return languages and length phases}\label{sec:phases55}
The existence of cofinite returns is not the only useful synchronization regime. In fact an exact-return language, when nonempty, has a rigid eventual form. Here and in this section a return has positive length.

\begin{lemma}[The return period]\label{lem:returnperiod55}
Suppose
\[
 \mathcal R=\{|w|>0:u_w=1\}\ne\varnothing,
 \qquad p=\gcd\mathcal R.
\]
There is an integer $g_p$ such that every multiple $dn$ with $n\ge g_p$ belongs to $\mathcal R$, and no length outside $p\N$ belongs to $\mathcal R$. In particular, every nonempty exact-return language is eventually one arithmetic progression.
\end{lemma}
\begin{proof}
Concatenation makes $\mathcal R$ additive. A finite subset of its elements already has gcd $p$: start with one length and successively decrease the gcd when necessary; a strictly decreasing chain of its positive divisors is finite. Divide these finitely many lengths by $p$. Their gcd is one, so the residue argument of Proposition~\ref{prop:return54} shows that their nonnegative sums contain every sufficiently large integer. Multiplying by $p$ proves the assertion. The exclusion of other lengths is the definition of the gcd.
\end{proof}

Fix a letter with product $a\in H$, and let $J$ be the closure of the products of words whose lengths are multiples of $p$, including the empty word.
\begin{lemma}[The phase subgroup]\label{lem:phasegroup55}
The group $J$ is an open normal subgroup of $H$. The quotient $H/J$ is cyclic of order dividing $p$, generated by $aJ$, and
\begin{equation}\label{eq:phasereachable55}
 \overline{\{u_w:|w|\equiv r\pmod p\}}=a^rJ
 \qquad(0\le r<p).
\end{equation}
If $H/H^\circ$ is finite, then $J^\circ=H^\circ$.
\end{lemma}
\begin{proof}
The compact closure defining $J$ is a group, by the positive-power argument, and is generated densely by the block alphabet $\mathcal A^p$. Also $a^p\in J$. If $x=u_w$ for a word $w$ of length $p$, then the executable word consisting of $p-1$ copies of the chosen letter, followed by $w$, followed by one more copy, has product $axa^{p-1}$ and length $2p$. Thus $axa^{p-1}\in J$, whence $axa^{-1}\in J$. It follows that $aJa^{-1}\subset J$. Applying this inclusion $p$ times and using $a^p\in J$ shows equality. Every letter $b$ has $ba^{p-1}\in J$, so $ba^{-1}\in J$. Thus all letters lie in $Ja$, normalize $J$, and their products of length $n$ lie in $Ja^n$. The finite union $\bigcup_{r=0}^{p-1}Ja^r$ is a closed group containing the letters, hence is $H$. This proves normality, finite index, and openness. Conversely, words consisting of $r$ copies of $a$ followed by arbitrary $p$-blocks have products dense in $Ja^r$, proving \eqref{eq:phasereachable55}. An open subgroup contains the identity component, and its own identity component is therefore the same one.
\end{proof}
The quotient order may be a proper divisor of $p$; the phase radii repeat whenever the physical cosets repeat. The phase index records executable length, not merely the physical component quotient.

Assume now that $H$ has finitely many components and put $K=H^\circ$. For $r\in\{0,\ldots,p-1\}$ define
\begin{equation}\label{eq:phaseradius55}
 \varepsilon_r=\max_{s,j,\ c\in H/K:\ c\subset a^rJ}
                    \rad_{Y_j}(F_{s,j}(c)).
\end{equation}
\begin{theorem}[The phasewise clocked threshold]\label{thm:phases55}
Under the sole synchronization assumption $\mathcal R\ne\varnothing$, for every $r$ and $\varepsilon\ge0$,
\[
 \sup_{N\equiv r\ (\mathrm{mod}\ p)}W_{N,\varepsilon}<\infty
             \quad\Longleftrightarrow\quad
             \varepsilon\ge\varepsilon_r.
\]
Equality is attained by a finite component-permutation realization on the indicated horizons. If $\varepsilon<\varepsilon_r$, then for each $k$ there is a constant $L_{k,r}$, independent of $N$, such that every error-$\varepsilon$ realization at $N\equiv r\pmod p$ satisfies
\[
                 \#\{1\le t\le N:|S_t|\le k\}\le L_{k,r}.
\]
The intervening widths are unrestricted. Consequently the full family is bounded exactly at and above $\max_r\varepsilon_r=\varepsilon_{\mathrm c}$; below it, divergence is asserted phase by phase, not necessarily along every horizon.
\end{theorem}
\begin{proof}
Every product of length $N\equiv r$ lies in $a^rJ$. Store its component and the seed and decode component centers as before. This proves the upper bound, including equality, with at most $|\mathcal S||H/K|$ labels. Centers on other components may be chosen arbitrarily when only this phase is under consideration.

For the converse fix a component $c\subset a^rJ$ and $s,j$ with radius greater than $\varepsilon$. We must count narrow cuts in every residue class, not only at block boundaries of one phase. For each $\ell\in\{0,\ldots,p-1\}$ choose $b_\ell\in\{0,\ldots,p-1\}$ congruent to $r-\ell$ modulo $p$. By \eqref{eq:phasereachable55} and the openness of components there is a fixed word $v_\ell$, of length $p_\ell\equiv\ell\pmod p$, with product in $a^{-b_\ell}c$. Indeed this component is contained in $a^\ell J$. Consider the block experiment on $J$ with alphabet $\mathcal A^p$ and target
\[
        G_\ell(h)=F_{s,j}(a^{b_\ell}h u_{v_\ell}),\qquad h\in J.
\]
It is a legal compact-convex interface. Its image on $K=J^\circ$ is $F_{s,j}(c)$, since $K$ is normal. Lemma~\ref{lem:returnperiod55} gives cofinite exact returns in block length for this alphabet.

Restrict the original machine to words that begin with $v_\ell$, end with $b_\ell$ copies of the chosen letter, and have $M=(N-p_\ell-b_\ell)/p$ arbitrary $p$-blocks in between. When $M\ge0$, absorb the fixed prefix into the initialization and the fixed suffix into the terminal decoder; stochastic averaging keeps that decoder in $Y_j$. This produces an error-$\varepsilon$ clocked block machine for $G_\ell$, with registers at original cuts $p_\ell+dm$, $0\le m\le M$. Theorem~\ref{thm:radius54} bounds its narrow cuts $1\le m\le M$ by a constant depending only on $k$ and the fixed data. The omitted cuts of residue $\ell$ number at most $p_\ell+b_\ell+1$: this also charges the block cut zero, while all counted cuts remain in $1,\ldots,N$. If $M<0$, the whole horizon is bounded by $p_\ell+b_\ell$. Sum these bounds over the $p$ residue classes. This counts all cuts $1,\ldots,N$, independently of all other register widths.
\end{proof}
The constants in this theorem depend on the fixed alphabet, return period and certificate, interface, error, width $k$, and chosen finite prefixes. They do not depend on the horizon or on the machine. The original constants $B,d,\chi,L_k$ in Section~\ref{sec:local54} use $d$ for a contraction defect. The return period here is denoted $p$ to keep the two parameters distinct.

\begin{proposition}[Different phases can have different thresholds]\label{prop:parity55}
There is an alphabet with return period two for which $W_{N,0}=1$ on every even horizon, but $W_{N,\varepsilon}\to\infty$ along odd horizons for every $0\le\varepsilon<\rho/2$.
\end{proposition}
\begin{proof}
Take $H=(\R/2\pi\Z)\times\Z/2\Z$ and letters $(\theta,1),(-\theta,1)$, with $\theta/2\pi$ irrational. The generated group is dense in $H$. An exact return must have zero signed count and hence even length; the two-letter return gives every positive even length. Here $J=(\R/2\pi\Z)\times\{0\}$. For one binary-mean query put $F(\phi,0)=0$ and $F(\phi,1)=\rho\cos\phi$, where $0<\rho\le1$. The even-phase target is identically zero and needs one label. The odd-phase radius in the norm $|\cdot|/2$ is $\rho/2$. Apply Theorem~\ref{thm:phases55}.
\end{proof}

\section{Infinite component groups and boundary attainment}\label{sec:profinite55}
This section retains independent finite-dimensional proofs of the strict-radius and exact translate-rank conclusions. Its output spaces are finite-dimensional. Let $H$ again be an arbitrary compact Hausdorff group, and set $K=H^\circ$. The quotient $H/K$ is profinite. We use two standard compact-group consequences of this fact and of Peter--Weyl theory \cite{Folland}: open normal subgroups form a neighborhood basis in a profinite group, and, for a continuous surjection $\pi:H\to G$ onto a compact Lie group, one has $\pi(K)=G^\circ$. To see the latter assertion, $\pi(K)$ is connected and normal, while $G/\pi(K)$ is a continuous group quotient of the profinite group $H/K$ and is therefore totally disconnected. This forces $G^\circ\subset\pi(K)$; the reverse inclusion is immediate. A continuous finite-dimensional representation of a profinite group has finite image, by the closed subgroup theorem and the no-small-subgroups property of a Lie group \cite{Hall}: an open normal subgroup can be placed inside a no-small-subgroups neighborhood, so its image is trivial.

Define the component radius
\begin{equation}\label{eq:R055}
                 R_0=\max_{s,j,g\in H}\rad_{Y_j}(F_{s,j}(gK)).
\end{equation}
Unlike $r(U)$, this quantity need not itself be realized by any finite quotient.

\begin{lemma}[Approximation by finite quotient radii]\label{lem:infradius55}
With $r(U)$ as in \eqref{eq:quotientradius55},
\[
                 \inf_{U\text{ open normal in }H}r(U)=R_0.
\]
\end{lemma}
\begin{proof}
Every open subgroup contains $K$, so $r(U)\ge R_0$. Fix $\eta>0$. Uniform continuity, for the finite family of interfaces, gives an identity neighborhood $V$ such that
\[
 \norm{F_{s,j}(xv)-F_{s,j}(x)}_j<\eta
       \quad(x\in H,\ v\in V,\ s,j).
\]
There is an open normal subgroup $U$ containing $K$ and contained in $KV$, by the profinite quotient basis. Write each $u\in U$ as $kv$ with $k\in K$, $v\in V$. Then $F_{s,j}(gu)$ lies within $\eta$ of $F_{s,j}(gk)$. A center for the latter component image consequently covers the former coset image with an additional error at most $\eta$. Thus $r(U)\le R_0+\eta$.
\end{proof}

\begin{lemma}[A legal finite-dimensional interface approximation]\label{lem:Lieapprox55}
For every $\eta>0$ there is a continuous homomorphism $\pi:H\to G\subset O(D)$ onto a compact Lie group and continuous maps $\widetilde F_{s,j}:G\to Y_j$ such that
\[
           \max_{s,j}\norm{F_{s,j}-\widetilde F_{s,j}\circ\pi}_\infty<\eta.
\]
Moreover, for every $g\in H$,
\[
 \left|\rad_{Y_j}(F_{s,j}(gK))-
 \rad_{Y_j}(\widetilde F_{s,j}(\pi(g)G^\circ))\right|<\eta.
\]
\end{lemma}
\begin{proof}
Approximate all coordinates of all $F_{s,j}$ by finitely many real representative functions. The direct sum of their orthogonal representations is $\pi$, and the approximants are continuous functions on its compact image $G$. Choose Euclidean inner products on the finitely many output spaces. Project each approximant to the closed convex set $Y_j$ by nearest-point projection. This projection is continuous and nonexpansive in the chosen Euclidean norm. Choose constants $a_j,b_j>0$ with $a_j\|v\|_j\le|v|_2\le b_j\|v\|_j$. An initial error less than $a_j\eta/b_j$ in $\|\cdot\|_j$ then gives a projected error less than $\eta$, by Euclidean nonexpansiveness. These constants are fixed before choosing the approximants. The closed subgroup $G\subset O(D)$ is a compact Lie group. Since $\pi(K)=G^\circ$, the two indicated images have Hausdorff distance less than $\eta$. Their constrained radii therefore differ by less than $\eta$.
\end{proof}

\begin{theorem}[Arbitrary compact groups on both strict sides]\label{thm:profinite55}
For any compact-group interface, $\varepsilon>R_0$ implies a finite stationary permutation realization. At $\varepsilon=R_0$, such a realization exists exactly when $r(U)\le R_0$ for some open normal $U$.

Suppose in addition that executable exact returns have all sufficiently large lengths. If $\varepsilon<R_0$, then for each $k$ there is a horizon-independent upper bound on the number of cuts of width at most $k$, with all intervening widths unrestricted. In particular $W_{N,\varepsilon}\to\infty$. No finite-component assumption is made in either assertion.
\end{theorem}
\begin{proof}
The stationary assertions follow from Lemma~\ref{lem:infradius55} and Theorem~\ref{thm:finitequotient55}. For the lower bound choose a component image of radius greater than $\varepsilon$ and then $\eta>0$ with $2\eta$ smaller than that strict difference. Apply Lemma~\ref{lem:Lieapprox55}. A machine correct for the original interface at error $\varepsilon$ is correct for the approximating interface at error at most $\varepsilon+\eta$. The chosen component radius of that interface is greater than $\varepsilon+\eta$. The projected commands still have cofinite exact returns and generate $G$ densely. Theorem~\ref{thm:radius54} therefore gives the desired occupation bound for the very same registers. This argument never assumes that $K$ is open or that executable products lying exactly in $K$ are dense.
\end{proof}
Theorem~\ref{thm:completeboundary56} now identifies this stationary boundary criterion with clocked boundedness, including $R_0>0$, and under the weaker approximate-return hypothesis. The zero-error boundary also admits the following independent translate-rank argument.

\begin{lemma}[Evaluation rank and translate dimension]\label{lem:evaluation56}
For a real function $f$ on a group, the supremum of the ranks of finite matrices $(f(y_lx_i))_{i,l}$ equals the dimension, possibly infinite, of the span of the functions $x\mapsto f(yx)$. These are left translates with the convention $(L_gf)(x)=f(g^{-1}x)$, after replacing $y$ by $g^{-1}$.
\end{lemma}
\begin{proof}
A finite-dimensional span bounds every evaluation rank. Conversely, if $f_1,\ldots,f_m$ are independent functions, their evaluation vectors $(f_1(x),\ldots,f_m(x))$ span $\R^m$: a proper span would have a nonzero annihilator, giving a linear dependence of the functions. Select $m$ evaluation vectors forming a basis. This gives an invertible evaluation matrix and proves the reverse inequality for every finite $m$.
\end{proof}

\begin{theorem}[The exact clocked boundary]\label{thm:zero55}
Assume cofinite executable exact returns. The following are equivalent:
\begin{enumerate}
\item $\sup_NW_{N,0}<\infty$;
\item all $F_{s,j}$ factor through one finite continuous quotient $H/U$, with $U$ open normal;
\item there is an exact stationary finite permutation realization.
\end{enumerate}
This includes groups with infinitely many connected components.
\end{theorem}
\begin{proof}
The implications from the second assertion to the third and from the third to the first are immediate. Suppose the first holds with bound $k$. Theorem~\ref{thm:profinite55} implies $R_0=0$, so every interface is constant on each $K$-coset.

Fix a real coordinate $f$ of one $F_{s,j}$. For any finite families of executable prefix products $x_i$ and suffix products $y_l$, use cofinite returns to pad all selected prefix words to one common length $A$ and all suffix words to one common length $B$. Choose an exact realization at horizon $A+B$ of width at most $k$. At its cut $A$ the matrix
\[
                            (f(y_lx_i))_{i,l}
\]
factors through the $k$ hidden labels: the left factors are the actual state distributions after the padded prefixes and the right factors are the actual conditional terminal expectations on the padded suffixes. Hence this matrix has rank at most $k$. Padding does not need to preserve hidden states; it only preserves the physical products used in the displayed entries. By density and continuity the same rank bound holds for arbitrary $x_i,y_l\in H$.

It follows that the span of the left translates of $f$ in $C(H,\R)$ has dimension at most $k$. Otherwise $k+1$ linearly independent translates would have an invertible evaluation matrix at some $k+1$ points, contrary to the preceding rank bound. Lemma~\ref{lem:evaluation56} formalizes this evaluation assertion. Equivalently it follows inductively: a nonzero function is nonzero somewhere, and a function outside the span of the preceding ones cannot have all its evaluations forced by theirs.

Take the sum $V$ of these translate spaces over all seed--query coordinates. It is finite-dimensional and invariant under left translation. Translation is continuous in the uniform norm and preserves that norm, so it gives a continuous representation of $H$ on $V$ with compact image; Haar averaging makes it orthogonal. Because $K$ is normal and every original coordinate is constant on $K$-cosets, $K$ acts trivially on every translate and thus on $V$. The representation factors through $H/K$. Its image is finite by the profinite finite-dimensional representation fact above. Its kernel $U$ is open normal. Since every original coordinate belongs to $V$, all interfaces are $U$-invariant and descend to $H/U$.
\end{proof}

\begin{proposition}[A nonattained zero-radius boundary]\label{prop:adic55}
On a compact group with infinitely many components, every positive error can have a bounded stationary realization while the zero-error clocked width is unbounded.
\end{proposition}
\begin{proof}
Let $H=\Z_2$ be the additive group of $2$-adic integers, with commands $0$ and $1$. Their positive products are dense, and the zero command supplies exact returns of every length. If $x=\sum_{n\ge0}x_n2^n$, $x_n\in\{0,1\}$, prescribe the binary success probability
\[
                         f(x)=\sum_{n\ge0}\frac{2x_n}{3^{n+1}}.
\]
Uniform convergence makes $f$ continuous. It is injective: at the first differing digit the contribution $2/3^{n+1}$ exceeds the sum $1/3^{n+1}$ of all possible later differences. Since $H$ is totally disconnected, $K=\{0\}$ and $R_0=0$. Truncation after $m$ digits factors through $\Z/2^m\Z$ and has uniform error at most $\sum_{n=m}^\infty2/3^{n+1}=3^{-m}$. Thus every positive tolerance admits a finite quotient machine. But the injective function $f$ cannot factor through any finite quotient. Theorem~\ref{thm:zero55} excludes bounded exact clocked width, and Theorem~\ref{thm:finitequotient55} excludes a finite exact stationary realization.
\end{proof}

\begin{theorem}[Linear exact width at a profinite boundary]\label{thm:adiclinear55}
For the interface of Proposition~\ref{prop:adic55},
\[
                         W_{N,0}=\Theta(N).
\]
More precisely, if $n$ is a power of two with $2n-1\le N$, then $W_{N,0}\ge n$, while $W_{N,0}\le N+1$. Every exact realization, with unrestricted intermediate widths, has at most $4k$ cuts in $\{1,\ldots,N\}$ of width at most $k$. Nevertheless, for each $\varepsilon>0$, $\sup_NW_{N,\varepsilon}<\infty$.
\end{theorem}
\begin{proof}
For nonnegative integers $z$ write $v_2(z+1)$ for the exponent of two in $z+1$. Incrementing the binary expansion, with $v_2(z+1)$ trailing ones, gives
\begin{equation}\label{eq:adicdifference55}
 b(z):=f(z+1)-f(z)=\frac53\,3^{-v_2(z+1)}-1.
\end{equation}
Let $n=2^m$, and form $B_n=(b(i+j))_{0\le i,j<n}$. Define a function on $\Z/n\Z$ by
\[
 \psi_m(x)=3^{-v_2(x)}\ (x\ne0),\qquad \psi_m(0)=3^{-m}.
\]
Here the valuation of a nonzero residue is independent of its representative. Since $1\le i+j+1\le2n-1$, its valuation equals that used in $\psi_m(i+j+1\bmod n)$, including the case $i+j+1=n$. Thus replacing row $i$ by the old row $(-i-1)\bmod n$ changes its entry at column $j$ to the kernel value at $j-i\bmod n$. This is the circulant with kernel $(5/3)\psi_m-1$.

On $\Z/n\Z$ one has the elementary identity
\[
            \psi_m=1-2\sum_{\ell=1}^m3^{-\ell}
                               \mathbf1_{2^\ell\mid x}.
\]
Use the eigenvector with entries $\exp(2\pi \mathrm i qj/n)$; its eigenvalue is the kernel sum against $\exp(2\pi \mathrm i qx/n)$. The sign is immaterial for subgroup indicators, whose transforms below are real. The Fourier eigenvalue at frequency zero is
\[
 \mu_0=\frac23 n6^{-m}>0.
\]
At frequency $q\ne0$ it is
\[
 \mu_q=-\frac{10n}{3}
           \sum_{\substack{1\le\ell\le m\\ n/2^\ell\mid q}}6^{-\ell}<0.
\]
Indeed a subgroup indicator has Fourier sum $n/2^\ell$ exactly when $n/2^\ell$ divides $q$, and zero otherwise. For $q\ne0$ the summand $\ell=m$ is always present. At $m=0$ the sole eigenvalue is $2/3$. Consequently $B_n$ is invertible for every dyadic $n$.

At an arbitrary command cut $t<N$, take a dyadic $n\le\min(t+1,N-t)$. For each $i<n$ choose a prefix of length $t$ with $i$ increment commands. For each $j<n$ choose two suffixes of length $N-t$ with respectively $j+1$ and $j$ increment commands. Fill all remaining positions with zero commands. The difference of their exact outputs is $b(i+j)$. Hence $B_n$ factors through the actual register $S_t$: one factor consists of the prefix state rows and the other of the differences of the two conditional suffix decoder columns. Invertibility gives $|S_t|\ge n$. The difference column need not be a legal decoder; it is used only for this rank bound on two legal experiments.

Taking $t=n-1$ proves $W_{N,0}\ge n$ whenever $2n-1\le N$. The largest dyadic $n\le(N+1)/2$ is greater than $(N+1)/4$, proving the linear lower order. Storing the number of increment commands uses at most $N+1$ labels and gives the exact upper bound.

For occupation, if $t<N$ and $|S_t|\le k$, the largest dyadic integer not exceeding $\min(t+1,N-t)$ is at most $k$, so $\min(t+1,N-t)<2k$. Such cuts lie within fewer than $2k$ positions of one of the two ends. Including the terminal cut gives at most $4k$ positive cuts. No width at another cut was restricted. The positive-error assertion is the finite-quotient truncation in Proposition~\ref{prop:adic55}.
\end{proof}
This supplies a sharp quantitative law on a genuinely infinite profinite group. Its linear exact-width obstruction and bounded positive-error realizations concern one fixed continuous interface; they are not a uniform approximation statement over all continuous outputs on that group.

\section{Binary, categorical, and observable interfaces}\label{sec:observable53}
\begin{corollary}[Component oscillation for binary means]\label{thm:classification53}
For a finite continuous binary-mean interface $f_{s,j}:H\to[-1,1]$, with the return hypothesis of Theorem~\ref{thm:radius54}, the exact bounded-width threshold in total variation is
\begin{equation}\label{eq:critical53}
 \epsc=\frac14\max_{s,j,c}\left(\max_{h\in c}f_{s,j}(h)-\min_{h\in c}f_{s,j}(h)\right).
\end{equation}
It is attained by a stationary permutation realization. Below it every fixed width has the occupation bound \eqref{eq:occupation54}.
\end{corollary}
\begin{proof}
In the norm $|\cdot|/2$, the radius of an interval is one quarter of its length. Its midpoint lies in $[-1,1]$. Apply Theorem~\ref{thm:radius54}.
\end{proof}
This recovers the scalar component theorem, including all of its nonuniform quantifiers. The vector theorem does not follow by applying this corollary separately to coordinates, because the coordinate centers need not belong to the output simplex.

\begin{proposition}[A three-outcome threshold]\label{prop:categorical54}
On the circle let $\omega=e^{2\pi i/3}$ and, for $j=0,1,2$, set
\[
 p_j(\phi)=\frac19\left|1+e^{i\phi}\omega^{-j}+e^{2i\phi}\omega^{-2j}\right|^2.
\]
For any dense circle command alphabet satisfying the return condition, the categorical law $p=(p_0,p_1,p_2)$ has critical total-variation error $2/3$. Each separately requested binary event $\{j\}$ has threshold $1/2$.
\end{proposition}
\begin{proof}
Orthogonality of the three characters gives $\sum_jp_j=1$; each entry is nonnegative. At $\phi=2\pi j/3$ the vector is the simplex vertex $e_j$. Hence every center $q\in\Delta_3$ has maximum distance at least
\[
 \max_j\TV(e_j,q)=1-\min_jq_j\ge2/3.
\]
The uniform center has distance at most $2/3$ from every point of $\Delta_3$, by convexity and the values at its vertices. Thus the constrained radius is exactly $2/3$. On the other hand each $p_j$ has range $[0,1]$, so its binary mean $2p_j-1$ has oscillation two. Apply Theorem~\ref{thm:radius54} and Corollary~\ref{thm:classification53}. In particular the categorical radius exceeds half the diameter, which here equals $1/2$.
\end{proof}

\begin{corollary}[Joint measure-once unitary laws]\label{cor:quantum54}
Let unitary commands have compact closure $H$ and satisfy the return condition. For a fixed unit seed $\psi_s$ and a finite measurement $E_{j,l}\ge0$, $\sum_lE_{j,l}=I$, prescribe the law
\[
 F_{s,j}(h)_l=\ip{h\psi_s}{E_{j,l}h\psi_s}.
\]
The critical error is its componentwise simplex radius in total variation. For closure $U(d)$ and one orthonormal-basis measurement, this radius is $1-1/d$.
\end{corollary}
\begin{proof}
The displayed map is a continuous simplex-valued polynomial in realified matrix entries. A closed unitary matrix group has finitely many components, so Theorem~\ref{thm:radius54} applies. For $U(d)$, every unit vector is reachable from the seed, and its squared coordinate magnitudes range over the whole simplex. The simplex vertices force radius at least $1-1/d$, attained by the uniform distribution.
\end{proof}
A query selects one complete measurement, not simultaneous answers to incompatible measurements. Several compatible finite outputs, with specified correlations, instead form one joint categorical law. Language recognition at a cutpoint is a different requirement.

For linear binary interfaces on a compact matrix group, write
\[
 f_{s,j}(h)=q_j^Thx_s\in[-1,1],\quad
 V_{\rm r}=\spanop\{hx_s:h\in H,s\in\mathcal S\},
\]
\[
 V_{\rm n}=\{x\in V_{\rm r}:q_j^Thx=0\text{ for all }h\in H,j\in\mathcal J\}.
\]
Both spaces are invariant. The observable space $V_{\rm o}=V_{\rm r}\cap V_{\rm n}^\perp$ is the orthogonal complement of $V_{\rm n}$ \emph{inside} $V_{\rm r}$, canonically realizing the observable quotient.

\begin{theorem}[Intrinsic zero-error criterion]\label{thm:observable53}
For these linear means, bounded exact clocked width is equivalent to finiteness of the generated group's image in $O(V_{\rm o})$, and to trivial action of $K$ on $V_{\rm o}$. An exact stationary permutation realization then exists.
\end{theorem}
\begin{proof}
Corollary~\ref{thm:classification53} identifies bounded exact width with constancy of all interface means on every component. Under that condition, for $k\in K$ and $g,h\in H$, normality gives
\[
 q_j^Th(k-I)gx_s=q_j^T(hkh^{-1})hgx_s-q_j^Thgx_s=0.
\]
Thus $(k-I)V_{\rm r}\subset V_{\rm n}$. On the invariant orthogonal complement $V_{\rm o}$ the difference also lies in $V_{\rm o}$, hence vanishes. The observable action factors through the finite component group. Conversely, if $K$ acts trivially on $V_{\rm o}$, the same difference is invisible and all means are componentwise constant. A finite image of the dense generated group has the same finite image closure from $H$, whose connected subgroup $K$ acts trivially. The zero-dimensional quotient is included.
\end{proof}

\begin{corollary}[The planar threshold]\label{cor:circle53}
Let an orthogonal planar alphabet contain $I$ and an irrational rotation, with any additional orthogonal letters. For signed coordinate seeds $\pm\rho e_1,\pm\rho e_2$, coordinate queries, and $0<\rho\le1$, the threshold is $\epsc=\rho/2$. At and above it one command label with a fair binary decoder suffices. Below it every fixed width has a finite occupation budget.
\end{corollary}
\begin{proof}
The identity component is $SO(2)$, and every specified mean has range $[-\rho,\rho]$ on each component of $H=SO(2)$ or $O(2)$. All midpoint decoders equal zero, so the seed/component register collapses to one label.
\end{proof}
For the noncommuting rational alphabet
\[
 I,\quad R=\begin{pmatrix}3/5&-4/5\\4/5&3/5\end{pmatrix},\quad
 F=\begin{pmatrix}1&0\\0&-1\end{pmatrix},
\]
one has $FRF=R^{-1}\ne R$. The rotation has infinite order: otherwise $(3+4i)/5$ and its inverse would be algebraic integers, whereas their rational sum $6/5$ is not an integer. This example has the sharp threshold above. We do not infer a badly approximable angle from rational matrix entries.

\begin{proposition}[Independent toral coordinates]\label{prop:torus53}
Suppose a component is identified with the full product torus $(\R/2\pi\Z)^m$, surjectively in independent coordinates, and a mean on it equals
\[
 a_0+\sum_{l=1}^m(a_l\cos\theta_l+b_l\sin\theta_l).
\]
Its contribution to \eqref{eq:critical53} is $\tfrac12\sum_l\sqrt{a_l^2+b_l^2}$.
\end{proposition}
\begin{proof}
Each summand has the symmetric range with amplitude $\sqrt{a_l^2+b_l^2}$, and all extrema are simultaneously attainable in the full product. Divide the resulting oscillation by four.
\end{proof}
On a proper subtorus, extrema must instead be taken subject to its relations. Ambient coordinates do not certify independence. Similarly, invisible directions do not create an obstruction: for $H=\{\diag(R_\theta,1)\}$, seed and query $e_3$ give constant mean one, whereas seed $\rho e_1$ and query $e_1$ have threshold $\rho/2$.

\begin{proposition}[Compact similarity]\label{prop:similarity53}
The classification applies to a finite alphabet in $GL(D,\R)$ whose generated group, including inverses, is uniformly bounded, with the executable return hypothesis retained. Conjugating to an invariant Euclidean metric changes neither output probabilities nor stochastic widths.
\end{proposition}
\begin{proof}
Uniform bounds on matrices and inverses make the closure a compact group. The matrix $Q=\int_Hh^Th\,d\mu(h)$ is positive definite and satisfies $A^TQA=Q$. Hence $Q^{1/2}AQ^{-1/2}$ is orthogonal. Transport the continuous output maps through this conjugacy; for a linear mean transport $x_s$ to $Q^{1/2}x_s$ and $q_j$ to $Q^{-1/2}q_j$. Exact identity products and their word lengths are preserved.
\end{proof}

\subsection{A curved categorical image with distinct component centers}
\begin{proposition}\label{prop:curved55}
Let $H=(\R/2\pi\Z)\times\Z/2\Z$, with an identity letter, an irrational rotation in the first factor, and a toggle of the second factor. Put
\[
 q_0=(\tfrac12,\tfrac13,\tfrac16),\qquad
 q_1=(\tfrac16,\tfrac13,\tfrac12),\qquad t=\tfrac1{12},
\]
and prescribe the three-outcome law
\[
 F(\theta,c)_l=q_{c,l}+t\cos(\theta-2\pi l/3),\qquad l=0,1,2.
\]
Each component image is a proper closed curve in the interior of the simplex, with unique constrained total-variation center $q_c$ and radius $t$. The union of both component images has radius $1/4$. Thus the sharp clocked threshold is $1/12$, and the stationary minimum at that threshold is exactly two labels.
\end{proposition}
\begin{proof}
The three cosines sum to zero and every coordinate is at least $1/12$, so $F$ is a legal law in the simplex interior. The first two independent trigonometric coordinates identify a nondegenerate ellipse in its affine plane. For a zero-sum vector $z\in\R^3$, one has $\frac12\norm{z}_1=\max_l|z_l|$. Each coordinate of the curve ranges over $[q_{c,l}-t,q_{c,l}+t]$. Hence for every center $y\in\Delta_3$,
\[
 \sup_\theta\TV(F(\theta,c),y)
       =t+\max_l|q_{c,l}-y_l|.
\]
This proves uniqueness of $q_c$ and radius $t$. For the union, the least possible value of $\max_{c,l}|q_{c,l}-y_l|$ is $1/6$, forced by the two endpoints in the first coordinate and attained by $y=(1/3,1/3,1/3)$. Its radius is therefore $t+1/6=1/4$. A one-label stationary machine has a constant decoder and cannot reach error $t$; two component labels attain it. The clocked threshold follows from Theorem~\ref{thm:radius54}.
\end{proof}

\begin{remark}
The total-variation diameter of a probability simplex is one, so half its diameter is $1/2$; the three-vertex example above has radius $2/3$. In the unitary corollary, $U(d)$ is connected. These facts explain respectively why coordinate midpoints do not solve the categorical problem and why the full-unitary component is the whole group.
\end{remark}

\section{Finite encodings and boundary quotients}\label{sec:effective54}
The general continuous theorem is not an algorithm from arbitrary real data. We now give an explicit finite input class for which its conclusions are effective. A rational orthogonal input consists of a finite alphabet in $O(D,\Q)$ containing $I$, finite seed and query sets, and rational polynomials $F_{s,j,l}$ in the matrix entries specifying categorical probabilities. The promise $F_{s,j}(H)\subset\Delta_{M_j}$ can itself be checked once the compact closure is computed. Real algebraic output numbers are encoded by polynomial equations and isolating data.

We use two established algorithmic inputs: computation of a finite defining description of the Zariski closure of a finitely generated rational matrix group \cite{DJK,NPSW}, and real quantifier elimination with effective semialgebraic connected components \cite{BPR}. The former is a terminating closure algorithm, not merely an enumeration of valid polynomial identities with an unknown stopping point. Compact orthogonal groups are real algebraic, so the real points of the computed closure in $O(D)$ give the Euclidean closure $H$ \cite{BJKP}. The complex algebraic components must not be substituted for the real connected components used here.

\begin{theorem}[Effective rational classification]\label{thm:effective54}
For a rational orthogonal input, the following are computable:
\begin{enumerate}
\item a semialgebraic description of every connected component of $H$, the finite group $C$, and the component of each command;
\item the critical total-variation error $\epsc$ as an exact real algebraic number, and an attaining stationary component-permutation machine with algebraic decoders;
\item for each algebraic $0\le\varepsilon<\epsc$ and integer $k\ge1$, an integer occupation constant $L_k$ valid for every error-$\varepsilon$ clocked realization;
\item for every fixed $N$ and algebraic $\varepsilon\ge0$, the least width $W_{N,\varepsilon}$ and an algebraic realization attaining it.
\end{enumerate}
These are terminating computability statements. No polynomial running-time bound, practical implementation of general closure/quantifier elimination, or finite-random-bit cost bound is asserted.
\end{theorem}
\begin{proof}
Compute $H$ using the cited closure algorithm, then its real connected components and algebraic sample points. The component containing $I$ is $K$. Products of sample points and semialgebraic membership tests determine the component multiplication table and the command images. In particular $K$ is now an effectively given compact semialgebraic group.

For a component $c$ and a query, the condition that $q$ be a radius-$t$ center is the real first-order formula
\begin{equation}\label{eq:QEcenter54}
 q\in\Delta_{M_j},\qquad
 \forall h\in c:\quad
 \sum_{l=1}^{M_j}|F_{s,j,l}(h)-q_l|\le2t.
\end{equation}
Absolute values have a finite semialgebraic description. Quantifier elimination gives the set of admissible $(q,t)$. Its least $t$ exists by compactness and is real algebraic; algebraic sampling supplies an attaining $q$. Maximize over the finite interface and components. This proves (1)--(2), including equality at the critical error.

We detail effectivity of the occupation proof, since computing the threshold alone would not supply it. Choose a component with radius greater than the given $\varepsilon$, and enumerate positive words until one lies in that component. Its rational product gives the prefix used in the proof of Theorem~\ref{thm:radius54}. The resulting map $f(h)=F_{s,j}(hu_v)$ is polynomial, so it can be used directly in place of $G$, with no approximation error.

For any finite tensor degree, all required polynomial functions are coefficients of the explicitly given tensor representation. Its fixed subspace is
\[
 \{z:\ \forall h\in K,
                 (R(h)-I)z=0\}.
\]
This linear subspace is effectively semialgebraic. Real algebraic sampling, successively outside the span of the vectors already chosen, gives an algebraic basis in finitely many steps. Its Gram matrix gives the orthogonal projection $P$ exactly. Thus every polynomial Haar moment is computable algebraically from $P\Phi(e)$; numerical integration of Haar measure is unnecessary.

Enumerate finite lists of squares of rational polynomials on $K$, discarding zero Haar integrals and normalizing the others. For a list $p_i$, compute $b_i=\int p_if\,d\mu$ and its constrained simplex radius by quantifier elimination. Rational polynomials are uniformly dense among real polynomial restrictions on a compact matrix group; the localization proof therefore guarantees that some enumerated list has radius strictly greater than $\varepsilon$. Strict algebraic comparison detects that list. Fix a positive rational $\Delta$ below the radius difference and a positive rational upper bound $L$ for all coefficient constants in Proposition~\ref{prop:residual54}. These give the certified residual $Q\ge\Delta/L$.

Next enumerate finite rational probability laws on executable words of a common length whose products lie in $K$, together with rational $d\in(0,1)$. The assertion \eqref{eq:gap54}, quantified over the finitely many unit vectors in the algebraic moving space, is semialgebraic: each maximum of finitely many inner products can be expressed by a finite disjunction. It is therefore decidable. Lemma~\ref{lem:gap54} and sufficiently small rational perturbations of its atom weights ensure that this search terminates. Padding by the identity letter makes every finite choice executable at one common length.

The resulting $B,d,L,\Delta,Q_0$ determine \eqref{eq:Lk54}. To avoid exact comparisons involving logarithms, take a rational $A'\ge\max(1,LQ_0/\Delta)$ and use
\[
 \frac{\log A}{-\log(1-d)}\le\frac{A'}d.
\]
An integer above $b+B-1+BA'/d$ is a computable occupation bound. This proves (3).

For (4), fix a candidate width $k$ and use $k$ labels at every cut, padding smaller registers with unused labels. Stochastic rows, initialization, and categorical decoder probabilities are finitely many variables in compact simplexes. For each of the finitely many length-$N$ words, the output vector is polynomial in these variables, and the target vector is rational. The total-variation inequalities are semialgebraic. Quantifier elimination therefore decides feasibility and supplies an algebraic point when feasible. A deterministic register storing the seed and entire prefix is a finite exact realization, so increasing $k$ must terminate. This proves the least-width claim for each finite horizon, without claiming a finite-horizon cutoff for all stationary or nonuniform realization questions.
\end{proof}
The distinction between an effective theorem and an implemented solver matters. The accompanying finite checks verify displayed exact examples and proof constants; they do not execute the general group-closure and quantifier-elimination construction. The finite-bit compiler for rational finite groups remains in the complete predecessor supplement with its separate hypotheses and costs. It is not used to reprice arbitrary algebraic decoders as cost-free finite-bit procedures.

\subsection{The least component-quotient realization}
The upper bound $|\mathcal S||C|$ can often be reduced. Here is an exact finite characterization of that reduction, with its scope stated explicitly. Put $X=\mathcal S\times C$. A \emph{component-quotient realization} has a deterministic state of the form $\pi(s,c)$, where $\pi:X\to Q$ is a surjection; its command update is induced by $(s,c)\mapsto(s,[u_a]c)$. It is stationary and has no additional state recording a representative word within a component.

Call a partition $\mathcal B$ of $X$ equivariant if
\[
 (s,c)\sim(s',c')\quad\Longrightarrow\quad
 (s,dc)\sim(s',dc')\qquad(d\in C).
\]
For a block $B\in\mathcal B$ and query $j$, set
\[
 A_{B,j}=\bigcup_{(s,c)\in B}F_{s,j}(c).
\]
\begin{theorem}[Finite quotient minimum]\label{thm:quotient54}
The least number of command labels in an error-$\varepsilon$ component-quotient realization is
\begin{equation}\label{eq:quotient54}
 \min\left\{|\mathcal B|:\mathcal B\text{ equivariant and }
         \rad_{Y_j}(A_{B,j})\le\varepsilon\text{ for every }B,j\right\},
\end{equation}
with value $+\infty$ if there is no admissible partition. For rational orthogonal categorical inputs this minimum and an attaining quotient are computable. If $\epsc=0$ and $\varepsilon=0$, the coarsest admissible partition is the future-output equivalence
\[
 (s,c)\sim(s',c')\quad\Longleftrightarrow\quad
 F_{s,j}(dc)=F_{s',j}(dc')\quad\text{for every }d\in C,j,
\]
where componentwise constant maps are identified with their values.
\end{theorem}
\begin{proof}
An equivariant partition induces well-defined command permutations on its blocks. Choosing a minimizing center for each $A_{B,j}$ gives exactly the stated error. Conversely, the fibers of $\pi$ in any component-quotient realization form an equivariant partition. Its decoder must approximate every executable product in each corresponding component. Such products are dense in that component, so continuity extends the error bound to $A_{B,j}$. Its decoder is therefore a permitted center. This proves \eqref{eq:quotient54}. The command generators suffice to check equivariance, since their images generate the finite group $C$.

There are finitely many partitions of $X$. For polynomial categorical data their radius tests are of the form \eqref{eq:QEcenter54}, with a finite disjunction over the relevant seed/component pairs, so the minimum is computable. In the exact componentwise constant case, equality of all future outputs is an equivariant equivalence relation. Every admissible partition refines it, and its own blocks have singleton output images for each query. It is therefore the unique coarsest admissible partition.
\end{proof}
The minimization in \eqref{eq:quotient54} is a deterministic quotient minimum, distinct from the unrestricted stationary minimum in Theorem~\ref{thm:stationaryminimum55} and from the unrestricted nonuniform minimum at the boundary. Randomized encodings can identify numerical behaviors without identifying component pairs. What is classified here is exactly the finite deterministic quotient construction that attains the radius threshold. Theorem~\ref{thm:effective54}(4) separately computes the unrestricted minimum at each fixed horizon.

\subsection{Stationary minimization by permutation closures}
\begin{theorem}[Effective unrestricted stationary minimum]\label{thm:effectivemin55}
For rational orthogonal commands generating $H$, rational polynomial categorical outputs, and algebraic $\varepsilon\ge0$, one can compute $M_\varepsilon$ and an attaining algebraic machine when it is finite. An identity letter is not required. At each proposed width $k$, at most $(k!)^{|\mathcal A|}$ augmented orthogonal group closures and finite real quantifier-elimination problems suffice.
\end{theorem}
\begin{proof}
First compute $H$, its real connected components, and their constrained radii as in Theorem~\ref{thm:effective54}(1)--(2); those steps do not use an identity letter. The group $H\subset O(D)$ has finitely many components. Theorem~\ref{thm:finitequotient55} shows that $M_\varepsilon=\infty$ below their maximum radius. Otherwise $M_\varepsilon\le |\mathcal S||H/H^\circ|$.

For $k$ between one and this upper bound enumerate all command-permutation tuples $\sigma$. The matrices
\[
                         \diag(u_a^{-1},\Pi_a)\in O(D+k,\Q)
\]
generate a compact closure which is exactly the joint group $\mathcal L_\sigma$ of Theorem~\ref{thm:stationaryminimum55}. Compute it by the same algebraic closure procedure. In \eqref{eq:permutationfeasibility55}, the unknown seed rows and decoder rows are constrained to simplexes. The mean output is bilinear in these rows, the group variable is restricted by computed defining equations whose real zero set inside the augmented orthogonal group is exactly $\mathcal L_\sigma$, not merely by some equations vanishing on it, and $h^{-1}=h^T$ in the physical block. Consequently all target coordinates are polynomial in the universally quantified matrix entries. Total variation is semialgebraic. Real quantifier elimination decides the resulting existential--universal formula and supplies algebraic rows whenever it is feasible. The first feasible $k$ is the minimum by purification. The finite upper bound proves termination of the entire search.
\end{proof}
This is a structural reduction of stationary minimization, not merely an enumeration of arbitrary stochastic transition parameters. It does not equate stationary and finite-horizon minima. Theorem~\ref{thm:stationarization56} identifies their limit under its additional return hypothesis. The two-cycle/three-cycle example in Proposition~\ref{prop:C655} already prevents replacing the permutation-feasibility test by deterministic component partitions.

\subsection{Algorithmic cost and arithmetic representation}
There are three distinct layers of effectivity. First, the group-closure procedure produces explicit defining equations; real-component decomposition and algebraic sampling then identify the component table. Second, radius computation and the stationary permutation tests are finite quantifier-elimination tasks on those equations. At fixed $k$, the stationary seed and output variables number $k(|\mathcal S|+\sum_jM_j)$ before normalization equalities are eliminated; the group variables can be taken among the $(D+k)^2$ augmented matrix entries. Thus the displayed finite search is independent of a horizon parameter. Its factorial number of discrete cases and the cost of the real-algebraic subproblems should not be mistaken for an efficient algorithm. No bound for the full closure-to-machine process in terms of only the original bit length is extracted here.

Third, occupation certificates in Theorem~\ref{thm:effective54}(3) add searches over localizer degrees and executable word laws. Their termination follows from strict analytic margins; the present proof supplies no a priori useful bound on the successful degree or word length. Fixed-horizon minimization, by contrast, has finitely many word constraints at each $N$ but their number is exponential in $N$ before repetitions are identified. These stages should be priced separately rather than hidden behind the word ``effective.''

All algebraic Haar moments in the localizer search are represented in a real algebraic number field with isolating data, using the fixed-subspace projector and its algebraic Gram inverse. Signs and strict comparisons are exact real-algebraic comparisons; a quadrature tolerance is not a certificate. Every rational command word has a rational matrix product and hence a rational target vector for rational polynomial readouts. The orthogonal closure fact used above is precisely that the Euclidean closure of a subgroup of $O(D)$ is the real zero set, inside $O(D)$, of its vanishing polynomial ideal, as used in \cite{BJKP}; replacing real connected components by complex algebraic ones would give the wrong quotient.

\section{Finite descriptions and decision complexity}\label{sec:complexity56}
The qualitative reduction is insensitive to the encoding of real numbers. Complexity statements are not. We first specify a finite input model and then give an explicit finite-precision counterpart to label width.

\begin{theorem}[Finite-group stationary minimization is $\exists\R$-complete]\label{thm:ETR56}
The input consists of a finite group $G$ given by its full multiplication table, finite seed and query sets, rational categorical vectors $F_{s,j}(g)$ for every $g\in G$, a nonnegative rational tolerance $\varepsilon$, and an integer $k\ge1$ in binary. Deciding whether $M_\varepsilon\le k$ is complete, under polynomial-time many-one reductions, for the existential theory of the reals. Hardness already holds for the trivial group, one query, and zero error. In particular the rational orthogonal polynomial-interface stationary problem is $\exists\R$-hard even with constant readouts and trivial dynamics.
\end{theorem}
\begin{proof}
Let $n=|G|$. If $k\ge n|\mathcal S|$, a deterministic seed--group register realizes every target exactly, so the answer is yes. Otherwise $k$ is bounded by a polynomial in the explicit input length, even though its encoding is binary.

For membership, Theorem~\ref{thm:equivariant56} permits permutation matrices $P_g$ forming a representation of $G$. Introduce real variables for their entries, constrain each entry $x$ by $x(x-1)=0$, impose every row and column sum equal to one, and impose
\[
                         P_1=I,
                 \qquad P_gP_h=P_{gh}\quad(g,h\in G).
\]
Introduce simplex variables $\alpha_s$ and categorical decoder rows $d_j(i)$. For every $s,j,g$ and output coordinate $c$, use a nonnegative auxiliary variable $z_{s,j,g,c}$ with
\[
 z_{s,j,g,c}\ge\pm\bigl((\alpha_sP_{g^{-1}}d_j)_c-F_{s,j}(g)_c\bigr),
 \qquad\sum_c z_{s,j,g,c}\le2\varepsilon.
\]
This is an existential real formula with rational coefficients, degree at most three, and polynomial size. Padding a representation by fixed unused labels handles an optimum strictly smaller than $k$. Theorem~\ref{thm:equivariant56} proves equivalence with the original unrestricted stationary problem. There is no group-closure computation in this explicit finite-group input model.

For hardness, let $A$ be a rational nonnegative matrix and let $k$ be a proposed nonnegative rank. Delete zero rows; the all-zero case is decided directly. Normalize each remaining row by its positive row sum to obtain a rational row-stochastic matrix $F$. This transformation is polynomial-time and preserves nonnegative rank. Interpret the rows as seeds and the columns as the categories of one query on the trivial group. A zero-error realization gives a stochastic factorization $F=\alpha d$ of inner dimension at most $k$, so the nonnegative rank is at most $k$.

Conversely, from any nonnegative factorization $F=UV$ delete a zero row of $V$ together with its column of $U$, and let $s_i=\sum_cV_{ic}>0$. Put $d_{ic}=V_{ic}/s_i$ and $\alpha_{bi}=U_{bi}s_i$. Then $d$ is row-stochastic and $\sum_i\alpha_{bi}=\sum_cF_{bc}=1$, so $\alpha$ is row-stochastic. Identity command rows yield an exact stationary realization. Thus $M_0$ equals the real nonnegative rank of $F$. Shitov's universality theorem, specifically its complexity consequence \cite[Theorem~2]{Shitov}, makes this decision problem $\exists\R$-complete. The last assertion follows by using the one-dimensional identity matrix as the only physical command and constant rational polynomial readouts.
\end{proof}
The upper bound is for an explicitly tabulated finite group. The hardness transfers to rational matrix input, but a polynomial-size existential description of an arbitrary compact matrix closure has not been supplied. Theorem~\ref{thm:effectivemin55} remains the separate termination result for that broader input class. The complexity obstruction above already occurs in randomized initialization and decoding; it is not caused solely by closure algorithms.

\subsection{A finite-table resource counterpart}
For a categorical interface, a dyadic table of precision $b$ has each probability in $2^{-b}\Z$. Besides peak label width $k$, charge its binary table length, all retained working registers, and the number of fair random bits used. A horizon-$N$ clocked implementation can store initialization, $N$ command tables, and terminal decoders. Its probability-table length is
\[
 O\left(b\left(|\mathcal S|k+N|\mathcal A|k^2+
                                k\sum_jM_j\right)\right).
\]
A uniform implementation is one finite algorithm that produces or accesses these tables from the encoded input, horizon, and accuracy. This requirement is additional to the original nonuniform model. An arbitrary real table, without an effective encoding, is not such an algorithm.

\begin{lemma}[Dyadic rounding of a probability row]\label{lem:dyadic56}
A probability row $p$ with $m$ entries has a dyadic row $\widehat p$ of precision $b$ such that
\[
                   \TV(p,\widehat p)\le(m-1)2^{-b}.
\]
One exact draw from $\widehat p$ uses at most $b$ independent fair bits.
\end{lemma}
\begin{proof}
Round the first $m-1$ entries down to multiples of $2^{-b}$ and put the remaining mass in the last entry. This is a probability row. Its total variation error equals the mass transferred to the last entry, which is at most $(m-1)2^{-b}$. For sampling, draw a uniform integer in $\{0,\ldots,2^b-1\}$ using $b$ bits and compare it with the integer cumulative row sums. Degenerate rows can use fewer bits.
\end{proof}

\begin{theorem}[Precision and random-bit budgets]\label{thm:bits56}
A categorical horizon-$N$ realization of peak width $k$ and error at most $\varepsilon$ admits a dyadic realization on the same registers whose error is at most
\begin{equation}\label{eq:clockround56}
       \varepsilon+\bigl((N+1)(k-1)+M_*-1\bigr)2^{-b},
                         \qquad M_* =\max_j M_j.
\end{equation}
Sampling initialization, all transitions, and one terminal answer requires at most $(N+2)b$ fair bits.

For a stationary compact-group realization of width $k$, the corresponding error bound can instead be made independent of $N$:
\begin{equation}\label{eq:stationaryround56}
                    \varepsilon+(k+M_*-2)2^{-b}.
\end{equation}
The stationary dyadic realization has at most $k$ labels, permutation command tables, table length
\[
 O\left(|\mathcal A|k\log(k+1)+
                   b\left(|\mathcal S|k+k\sum_jM_j\right)\right),
\]
and uses at most $2b$ fair bits in total for one experiment of any length. Given exact algebraic encodings of the selected initialization and decoder rows, the rounding and sampling are effective finite procedures.
\end{theorem}
\begin{proof}
Apply Lemma~\ref{lem:dyadic56} to every row. Couple the initial draws maximally and then, as long as the two labels agree, maximally couple their next transition draws. The probability of any disagreement is at most the sum of the initialization and $N$ row errors, hence at most $(N+1)(k-1)2^{-b}$. Conditioned on matching terminal labels, the terminal categorical rows have total variation error at most $(M_*-1)2^{-b}$. The coupling inequality gives \eqref{eq:clockround56}, uniformly over words and queries. The sample count is the number of draws times $b$.

For the stationary claim, first use Theorem~\ref{thm:purification55}. Its transition matrices are permutations and are already dyadic at every precision. Round only initialization and terminal categorical rows. The initialization discrepancy cannot increase under a permutation, so the same coupling gives \eqref{eq:stationaryround56} with no factor $N$. There are no random transition draws, leaving only initialization and one terminal answer. A permutation is stored as $k$ integer images, giving the displayed description bound. Exact comparison of algebraic numbers with dyadic rationals computes the necessary floors and handles equality, after which integer cumulative sampling is immediate.
\end{proof}
For a prescribed extra error $\eta>0$, choosing $b$ with the added term in \eqref{eq:clockround56} or \eqref{eq:stationaryround56} at most $\eta$ gives an explicit precision budget. The persistent label uses $\lceil\log_2k\rceil$ bits; an implementation of the sampler also retains a $b$-bit random integer and bounded indexing registers while sampling. These temporary registers are not disguised as free persistent randomness. The description and random-bit bounds do not assert a polynomial time bound for constructing an algebraic optimum. They also do not turn an exact irrational realization into an exact finite-bit one: the extra tolerance is essential in that case.

\section{Classical structure and the additional quantifiers}
Classical stochastic-semigroup and nonnegative-matrix-group structure,
including Flor \cite{Flor} and Schwarz \cite{Schwarz58}, supplies recurrent
idempotents and permutation actions on their finite recurrent simplexes.
The purification proof uses that structure inside the joint physical and
stochastic closure. The extra conclusion keeps the prescribed closed
numerical error balls and the original label budget, then descends the
hidden action to the physical group.

Fijalkow--Paperman \cite[Theorems~1 and 10]{FP} characterize advice
regular languages and prove neutral-letter collapse. In fact their
Myhill--Nerode proof preserves the relevant deterministic bound $K$;
we do not claim that preservation of a state bound, by itself, is new.
The difference here is the stochastic simplex, common closed numerical
error, approximate executable returns and an occupation obstruction with
unrestricted intermediate widths. Our proof keeps the homogeneous
neutral table inside the terminal decoder rather than identifying it
with a transition identity.

Rational and weighted-series minimization \cite{BR,Paz} concerns linear
representations and Hankel spaces. Positive realization \cite{BF04}
adds a cone constraint. Neither linear dimension nor a separate
nonnegative factorization at each cut is used here as a substitute for
one compatible stochastic realization. Shitov's exact complexity
statement \cite[Theorem~2]{Shitov} is imported in the finite-group
hardness proof, not re-proved as a realization theorem.

For group and quantum automata, Brodsky--Pippenger
\cite[Theorems~3.3 and 3.6]{BP} compare cut-point languages, whereas our
physical action retains a full continuous numerical interface and its
closed error tolerance. Distribution metrics \cite{FSZ} and topological
recognition \cite{Steinberg} give adjacent equivalence and dynamical
frameworks; they are not assumed to identify the same-width stochastic
minimum. Classical hidden Markov realization \cite{Heller58} imposes
process consistency that is not requested by a single terminal query.
No online or general irreversible-group-quotient conclusion follows.

The effectivity assertions must be kept separate. A finite continuous
quotient may exist without an input encoding. Algebraic closure
algorithms permit computation for rational matrix input. An explicitly
tabulated finite group has the displayed existential-real decision
classification. Constructing a minimum realization and compiling a
supplied encoded realization are further tasks with different costs.
The finite-precision output sizes are not, in general, runtime estimates.
This theorem-level comparison is author-side positioning, not independent
priority clearance.

\subsection{Controlled stochastic processes}
Singh, James and Rudary \cite[Theorems~2 and~6]{SJR2004} distinguish
finite hidden-state dimension from linear predictive dimension. Their
first theorem bounds the rank of the system-dynamics matrix by the
number of POMDP states; their second characterizes finite linear
predictive representations, whose update coefficients need not be
nonnegative. Neither rank counts distinct future rows.
Our causal theorem instead uses repeated nondisturbing probes to make
finitely many selected future responses statistically distinguishable
by one common experiment. It then tests every horizon-dependent
nonnegative instrument across the same register. The new condition is
uniform total variation of whole adaptive transcripts. It is not
inferred from a linear-rank calculation and does not claim a general
HMM classification. The proof supplies the stated result, not an
independent certification of priority across the realization literature.

\begin{thebibliography}{99}
\item[]\textit{Stochastic semigroups and positive realization}
\bibitem{Flor} P. Flor, On groups of non-negative matrices, \emph{Compositio Mathematica} \textbf{21} (1969), no. 4, 376--382.
\bibitem{Schwarz58} \v S. Schwarz, On the structure of the semigroup of stochastic matrices, \emph{Publications of the Mathematical Institute of the Hungarian Academy of Sciences, Series A} \textbf{9} (1964), 297--311.
\bibitem{BF04} L. Benvenuti and L. Farina, A tutorial on the positive realization problem, \emph{IEEE Transactions on Automatic Control} \textbf{49} (2004), no. 5, 651--664.
\bibitem{Heller58} A. Heller, On stochastic processes derived from Markov chains, \emph{Annals of Mathematical Statistics} \textbf{36} (1965), no. 4, 1286--1291.
\item[]\textit{Automata, weighted series, and recognition}
\bibitem{Rabin} M. O. Rabin, Probabilistic automata, \emph{Information and Control} \textbf{6} (1963), no. 3, 230--245.
\bibitem{Paz} A. Paz, \emph{Introduction to Probabilistic Automata}, Academic Press, New York, 1971.
\bibitem{BR} J. Berstel and C. Reutenauer, \emph{Noncommutative Rational Series with Applications}, Encyclopedia of Mathematics and its Applications 137, Cambridge University Press, Cambridge, 2011.
\bibitem{FP} N. Fijalkow and C. Paperman, Monadic second-order logic with arbitrary monadic predicates, arXiv:1709.03117 (2017).
\bibitem{BP} A. Brodsky and N. Pippenger, Characterizations of 1-way quantum finite automata, \emph{SIAM Journal on Computing} \textbf{31} (2002), no. 5, 1456--1478.
\bibitem{AF57} A. Ambainis and R. Freivalds, 1-way quantum finite automata: strengths, weaknesses and generalizations, arXiv:quant-ph/9802062 (1998).
\bibitem{FSZ} Y. Feng, L. Song, and L. Zhang, Distribution-based bisimulation and bisimulation metric in probabilistic automata, arXiv:1512.05027 (2015).
\bibitem{Steinberg} B. Steinberg, Topological dynamics and recognition of languages, arXiv:1306.1468 (2013).
\item[]\textit{Compact groups and orbit geometry}
\bibitem{Folland} G. B. Folland, \emph{A Course in Abstract Harmonic Analysis}, second ed., Chapman \& Hall/CRC, Boca Raton, 2016.
\bibitem{Hall} B. C. Hall, \emph{Lie Groups, Lie Algebras, and Representations: An Elementary Introduction}, second ed., Graduate Texts in Mathematics 222, Springer, Cham, 2015.
\bibitem{Szarek57} S. J. Szarek, Metric entropy of homogeneous spaces, \emph{Banach Center Publications} \textbf{43} (1998), 395--410.
\bibitem{BG2012} J. Bourgain and A. Gamburd, A spectral gap theorem in $SU(d)$, \emph{Journal of the European Mathematical Society} \textbf{14} (2012), no. 5, 1455--1511.
\bibitem{BreuillardGelander57} E. Breuillard and T. Gelander, On dense free subgroups of Lie groups, \emph{Journal of Algebra} \textbf{261} (2003), no. 2, 448--467.
\bibitem{BreuillardGelanderCorrection57} E. Breuillard and T. Gelander, On dense free subgroups of Lie groups---revisited, arXiv:2605.20568 (2026).
\item[]\textit{Algebraic computation and complexity}
\bibitem{BPR} S. Basu, R. Pollack, and M.-F. Roy, \emph{Algorithms in Real Algebraic Geometry}, second ed., Algorithms and Computation in Mathematics 10, Springer, Berlin, 2006.
\bibitem{BJKP} V. D. Blondel, E. Jeandel, P. Koiran, and N. Portier, Decidable and undecidable problems about quantum automata, \emph{SIAM Journal on Computing} \textbf{34} (2005), no. 6, 1464--1473.
\bibitem{DJK} H. Derksen, E. Jeandel, and P. Koiran, Quantum automata and algebraic groups, \emph{Journal of Symbolic Computation} \textbf{39} (2005), nos. 3--4, 357--371.
\bibitem{deGraaf} W. A. de Graaf, Computing the Zariski closure of a finitely generated matrix group, arXiv:2505.02577v2 (2026).
\bibitem{NPSW} K. Nosan, A. Pouly, S. Schmitz, M. Shirmohammadi, and J. Worrell, On the computation of the Zariski closure of finitely generated groups of matrices, arXiv:2106.01853v6 (2025).
\bibitem{Shitov} Y. Shitov, A universality theorem for nonnegative matrix factorizations, arXiv:1606.09068v2 (2018).
\bibitem{Ramsey} F. P. Ramsey, On a problem of formal logic, \emph{Proceedings of the London Mathematical Society} (2) \textbf{30} (1930), no. 1, 264--286.
\bibitem{SJR2004}
S.~Singh, M.~R.~James, and M.~Rudary,
\emph{Predictive state representations: A new theory for modeling dynamical systems},
Proceedings of the Twentieth Conference on Uncertainty in Artificial Intelligence,
AUAI Press, 2004, 512--519. See Theorems~2 and~6; arXiv:1207.4167.
\end{thebibliography}

\end{document}
